\documentclass[10pt,a4paper]{article}

% ----------------- Paquetes -------------------%
\usepackage[T1]{fontenc}
\usepackage[utf8]{inputenc}
\usepackage[a4paper,margin=2.5cm]{geometry}
\usepackage{mathptmx}             
\usepackage{changepage} 
\usepackage{setspace}
\usepackage[spanish]{babel}
\usepackage[labelfont=bf,labelsep=space]{caption}
\usepackage{amssymb}
\usepackage{amsmath}
\usepackage{ebgaramond}
\usepackage{placeins}
\usepackage{etoolbox}
\usepackage{setspace}
\usepackage{parskip} 
\usepackage{tcolorbox}
\usepackage{needspace}
\usepackage{xparse}
\usepackage{ocgx2}
\usepackage{hyperref}
\usepackage{hypcap}
\usepackage{soul}\sethlcolor{yellow}% resaltado amarillo: \hl{texto} (Panel TCEE, ⌃H)


%%%%%%%%%%%%%%%%%%%%%%%%%%%%%%%%%%%%%%%%%%%%%%%%%%%%%%%%%%%%%%%%
% --- Fija primero tus valores globales "buenos" ---
\setstretch{1.2}                 % Interlineado global
\setlength{\parskip}{0.7\baselineskip}
\setlength{\parindent}{0pt}
\newlength{\GlobalParskip}
\setlength{\GlobalParskip}{\parskip}

\newcounter{eqnum}[subsection]
\renewcommand{\theeqnum}{\arabic{eqnum}}
\newcounter{alphaitem}
\newcounter{numitem}

%% ------------------- Recuadros----------------------%%
\newcommand{\puntosclave}[1]{%
  \begin{tcolorbox}[
    colframe=black,
    title={Puntos clave},
    before upper={\setlength{\parindent}{0.5cm}\setlength{\parskip}{0.5\baselineskip}}
  ]
  #1
  \end{tcolorbox}
}

\newcommand{\modificaciones}[1]{
  \begin{tcolorbox}[
    colback=white,
    colframe=red,
    title={Modificaciones a realizar},
    before upper={\setlength{\parindent}{0.5cm}\setlength{\parskip}{0.5\baselineskip}}
  ]
  #1
  \end{tcolorbox}
}

%% ------------------- Preguntas TEST----------------------%%
% Opciones: radio (single-choice)
\NewDocumentCommand{\OptRadio}{m m m}{%
  \noindent\CheckBox[radio,name=#1,value=#2,width=1.2ex,height=1.2ex]{}%
  \;\textbf{#2.}~#3\par
}

% Opciones: checkbox (multi-choice)
\NewDocumentCommand{\OptCheck}{m m}{%
  \noindent\CheckBox[name=#1,width=1.2ex,height=1.2ex,bordercolor={0 0 0}]{}%
  \;#2\par
}

% Pregunta single-choice (radio)
% Args: 1=id, 2=enunciado, 3=opciones, 4=solución+justificación, 5=imagen (o {}), 6=origen
\NewDocumentCommand{\PreguntaTestSingle}{m m m m m m}{%
  \par\medskip
  \noindent\textbf{#1.}~#2\par\smallskip

    % Imagen (si no es {})
    \IfBlankTF{#5}{}{%
      \begin{center}
        \includegraphics[width=0.75\linewidth]{#5}
      \end{center}
    }

    \begin{Form}
      #3
    \end{Form}

    \par\smallskip
    \switchocg{sol-#1}{\fbox{Mostrar / ocultar solución}}%
    \hfill{\footnotesize #6}\par\smallskip

    \begin{ocg}{Solución}{sol-#1}{0}
      \noindent #4
    \end{ocg}
  \par\medskip
}

% Pregunta multi-choice (checkboxes)
\NewDocumentCommand{\PreguntaTestMulti}{m m m m m m}{%
  \par\medskip
  \noindent\textbf{#1.}~#2\par\smallskip

    % Imagen (si no es {})
    \IfBlankTF{#5}{}{%
      \begin{center}
        \includegraphics[width=0.75\linewidth]{#5}
      \end{center}
    }
    \begin{Form}
      #3
    \end{Form}

    \par\smallskip
    \switchocg{sol-#1}{\fbox{Mostrar / ocultar solución}}%
    \hfill{\footnotesize #6}\par\smallskip

    \begin{ocg}{Solución}{sol-#1}{0}
      \noindent #4
    \end{ocg}
  \par\medskip
}

%% ------------------- Listas----------------------%%
    % --- Lista alfabética ---
    \newenvironment{la}
      {\begin{list}{\alph{alphaitem})}{%
          \usecounter{alphaitem}%
          \setlength{\leftmargin}{1cm}%
          \setlength{\parskip}{3pt}% 
          \setlength{\itemsep}{0pt}% ← espacio entre ítems
          \setlength{\parsep}{0pt}%  ← párrafos dentro de un ítem
      }}
      {\end{list}}
      
    % --- Lista numérica ---
    \newenvironment{lnum}
      {\begin{list}{\arabic{numitem}.}{%
          \usecounter{numitem}%
          \setlength{\leftmargin}{1cm}%
          \setlength{\parskip}{3pt}% 
          \setlength{\itemsep}{0pt}% ← espacio entre ítems
          \setlength{\parsep}{0pt}%  ← párrafos dentro de un ítem
      }}
      {\end{list}}
      
% ------------------- Imágenes----------------------%
    \graphicspath{{Img/}}% Rutas por defecto 
    % --- Imagen adaptativa ---
    % Registros de dimensión definidos una sola vez
    \newdimen\imgcap
    \newdimen\imgmin
    \newdimen\imgmax
    \newdimen\imgremain
    \newdimen\imgh
    \newdimen\imgneed
    
    \newcommand{\imagenfit}[3]{%
      \addvspace{10pt}% separación antes
      \par\begingroup
        % Parámetros de composición (asignaciones, no \newdimen)
        \imgcap=1\baselineskip            % alto estimado de título+fuente
        \imgmin=0.15\textheight
        \imgmax=0.15\textheight
        \imgremain=\pagegoal
        \ifdim\pagegoal=\maxdimen \imgremain=0pt\fi
        \advance\imgremain by -\pagetotal
        \advance\imgremain by -\imgcap
        % Decisión de altura destino
        \ifdim\imgremain<\imgmin
          \newpage
          \imgh=0.20\textheight
        \else
          \imgh=\imgremain
          \ifdim\imgh>\imgmax \imgh=\imgmax \fi
        \fi
        % Reservar espacio para evitar cortes del bloque completo
        \imgneed=\imgh \advance\imgneed by \imgcap
        \Needspace{\imgneed}
        % Cuerpo
        \noindent\begin{minipage}{\textwidth}
          \centering
          \captionof{figure}{\textemdash\ #2}\par\vspace{0.3em}% título arriba
          \includegraphics[width=\linewidth,height=\imgh,keepaspectratio]{\detokenize{#1}}\par
          \vspace{0.3em}\small\textit{Fuente: }#3% fuente abajo
        \end{minipage}%
      \endgroup
      \par\addvspace{10pt}%
    }

%------------ Bloque de RECUADRO ESPECIAL -------------%
    \newenvironment{specialblock}[1]{%
      \begin{quote} % crea sangría a izquierda y derecha
      \small % texto más pequeño
      \noindent\textbf{#1}\\[-0.3em] % título en negrita
      \rule{\linewidth}{0.4pt}\\[0.6em]}{
      \end{quote}}
      
%-------------------------- Portada -----------------------%
\newcommand{\oldtitle}[1]{\def\OldTitle{#1}}
\newcommand{\changerationale}[1]{\def\ChangeWhy{#1}}

\newcommand{\changebox}{%
\begin{tcolorbox}[title={Historial de cambios del tema}]
\textbf{Título anterior:} \OldTitle\par
\textbf{Motivo del cambio:} \ChangeWhy
\end{tcolorbox}
}

%-------------------------- Author Font -----------------------%
    \newcommand{\authorfont}[1]{{\fontfamily{EBGaramond-TLF}\selectfont #1}}

%-------------------------- Anexos -----------------------%
    \makeatletter
    \newcounter{anexo}
    \renewcommand{\theanexo}{\Roman{anexo}}
    \newcommand{\anexo}[1]{%
      \clearpage
      \phantomsection
      \refstepcounter{anexo}%
      \noindent\textbf{Anexo \theanexo.- }#1\par\medskip
      \addcontentsline{stc}{section}{Anexo \theanexo.- #1}%
    }
    \makeatother

    %-------------------------- Lista de figuras (personal) -----------------------%
\makeatletter
\newcommand{\MyListOfFigures}{%
  \clearpage
  \phantomsection
  {\centering\bfseries\Large Índice de figuras\par}%
  \medskip
  % Entrada en nuestro índice de contenido (stc)
  \addcontentsline{stc}{section}{Índice de figuras}%
  % Recuperación del listado (sin cabecera propia)
  \begingroup
    \parindent=0pt\parskip=2pt
    \@starttoc{lof}%
  \endgroup
}
\makeatother

%-------------------------- Bibliografía -----------------------%
\makeatletter
% Contador para las referencias
\newcounter{myref}
% Cabecera de la bibliografía, entra en el índice 'stc'
\newcommand{\MyBibliography}{%
  \clearpage
  \phantomsection
  {\centering\bfseries\Large Bibliografía y referencias\par}%
  \medskip
  \addcontentsline{stc}{section}{Bibliografía y referencias}%
  \setcounter{myref}{0}%
}
\newcommand{\MyBibItem}[4]{%
  \refstepcounter{myref}%
  \noindent[\themyref]\ #1.\ \emph{#2}. #3. #4\par
}
\makeatother

%--------- Establecer cuatro niveles de índice --------%
    \setcounter{tocdepth}{4}   
    \setcounter{secnumdepth}{4}% numera hasta \paragraph

%%%%%%%%%%%%%%%%%%%%%%%%%%%%%%%%

\makeatletter

    %--------- Interlineado Global --------%
    \edef\GlobalBaseLineStretch{\baselinestretch}
    
    %--------- Espaciado de seccinones --------%
    \renewcommand\section{\@startsection{section}
      {1}{\z@}
      {2.5ex \@plus 1ex \@minus .2ex} % espacio antes
      {1.5ex \@plus .2ex}             % espacio después
      {\normalfont\Large\bfseries}    % formato del título
    }
    \renewcommand\subsection{\@startsection{subsection}{2}{\z@}
      {3ex \@plus 0.8ex \@minus .2ex}
      {1.5ex \@plus .2ex}
      {\normalfont\large\bfseries}
    }
    \renewcommand\subsubsection{\@startsection{subsubsection}{3}{\z@}
      {3ex \@plus 0.5ex \@minus .2ex}
      {1ex \@plus .2ex}
      {\normalfont\normalsize\bfseries}
    }
    \renewcommand\paragraph{\@startsection{paragraph}{4}{\z@}%
    {-2ex \@plus -0.5ex \@minus -.2ex}% espacio antes
    {0.8ex \@plus .2ex}% espacio después
    {\normalfont\normalsize\itshape}}% ← cursiva en lugar de negrita

    %--------- Ecuaciones --------%   
    % Variables globales para almacenar títulos y notas
    \gdef\leq@current@section{}%
    \gdef\leq@current@subsection{}%
    \gdef\leq@last@printed@section{}%
    \gdef\leq@last@printed@subsection{}%
    
    % Comando para añadir nota (invisible en el cuerpo)
    \newcommand{\eqnote}[1]{%
      \addcontentsline{leq}{leqnote}{#1}%
    }
    
    % Comando para ecuaciones
    \newcommand{\eqblock}[2]{%
      % Verificar si necesitamos imprimir el título de sección
      \ifx\leq@current@section\leq@last@printed@section\else
        \addcontentsline{leq}{leqsection}{\leq@current@section}%
        \global\let\leq@last@printed@section\leq@current@section
        \gdef\leq@last@printed@subsection{}% Reset subsección al cambiar sección
      \fi
      % Verificar si necesitamos imprimir el título de subsección
      \ifx\leq@current@subsection\leq@last@printed@subsection\else
        \ifx\leq@current@subsection\@empty\else
          \addcontentsline{leq}{leqsubsection}{\leq@current@subsection}%
          \global\let\leq@last@printed@subsection\leq@current@subsection
        \fi
      \fi
      % Ecuación
      \refstepcounter{eqnum}%
      \begin{equation*}
        #1\tag{E.\theeqnum}
      \end{equation*}%
      \addcontentsline{leq}{leqt}{[E.\theeqnum] -- #2}%
      \addcontentsline{leq}{leqm}{#1}%
    }
    
    %%% ----- Índice de ecuaciones ----------%%%
    % Tamaño para []
    \newcommand{\leqIndexFont}{\normalsize}
    \newcommand{\leqTitleFont}{\tiny}
    \newcommand{\leqNoteFont}{\tiny}
    % Cabecera de sección 
    \newcommand*\l@leqsection[2]{%
      \addvspace{25pt}% Mayor separación antes de nueva sección
      \noindent\leqIndexFont
      \makebox[\linewidth]{\rule{.38\linewidth}{0.4pt}\ \ \textbf{  \MakeUppercase{#1}}\ \ \rule{.38\linewidth}{0.4pt}}%
      \par\vspace{-10pt}
    }
    % Cabecera de subsección 
    \newcommand*\l@leqsubsection[2]{%
      \par\addvspace{15pt}%
      \noindent\leqIndexFont #1
      \par\vspace{-7pt}
    }
    % Título de ecuación 
    \newcommand*\l@leqt[2]{%
    \par\addvspace{10pt}%
      \par\noindent\leqTitleFont #1\par
    }
    % Miniatura de ecuación
    \newcommand*\l@leqm[2]{%
      \vspace{0.3\baselineskip}%
      \par\scriptsize\noindent\hspace*{1.5em}\(#1\)\par%
    }
    % Nota de ecuación 
    \newcommand*\l@leqnote[2]{%
      \vspace{0.2\baselineskip}%
      \par\noindent\leqNoteFont\hspace*{1.5em}\textbf{Nota.--} #1\par\vspace{0.5\baselineskip}%
    }
    % Generación del índice
    \newcommand{\mylistofequations}{%
      \clearpage
      \phantomsection
      % Título visual de la lista (ajústalo a tu gusto):
      {\centering\bfseries\Large Índice de ecuaciones\par}
      % Entrada en nuestro índice 'stc':
      \addcontentsline{stc}{section}{Índice de ecuaciones}%
      % Llamada a tu lista real (usa el nombre exacto de tu macro):
      \@starttoc{leq}%
    }
    
    %--------- Captura de títulos de sección --------%
    \let\leq@original@section\section
    \let\leq@original@subsection\subsection
    % Flag para saber si estamos en una sección asterisco
    \newif\ifLeqStarSection
    \LeqStarSectionfalse
    
    % Redefinir \section CON CONTROL DE ASTERISCO
    \renewcommand{\section}{%
      \@ifstar{\leq@section@star}{\@dblarg\leq@section@nostar}%
    }
    \def\leq@section@star#1{%
      \global\LeqStarSectiontrue
      \gdef\leq@current@section{#1}%
      \gdef\leq@current@subsection{}%
      \leq@original@section*{#1}%
      \global\LeqStarSectionfalse
    }
    \def\leq@section@nostar[#1]#2{%
      \global\LeqStarSectionfalse
      \gdef\leq@current@section{#2}%
      \gdef\leq@current@subsection{}%
      \leq@original@section[#1]{#2}%
    }
    % Redefinir \subsection
    \renewcommand{\subsection}{%
      \@ifstar{\leq@subsection@star}{\@dblarg\leq@subsection@nostar}%
    }
    \def\leq@subsection@star#1{%
      \global\LeqStarSectiontrue
      \gdef\leq@current@subsection{#1}%
      \leq@original@subsection*{#1}%
      \global\LeqStarSectionfalse
    }
    \def\leq@subsection@nostar[#1]#2{%
      \global\LeqStarSectionfalse
      \gdef\leq@current@subsection{#2}%
      \leq@original@subsection[#1]{#2}%
    }

    % ----- Restauración de valores Globales de estilo ----------%
    % Flags
    \newif\ifInFootnote
    \InFootnotefalse
    \pretocmd{\@footnotetext}{\InFootnotetrue}{}{}
    \apptocmd{\@footnotetext}{\InFootnotefalse}{}{}
    % Profundidad de listas (soporta anidadas)
    \newcount\MyListDepth \MyListDepth=0
    \AtBeginEnvironment{la}{\advance\MyListDepth by 1}
    \AfterEndEnvironment{la}{\advance\MyListDepth by -1}
    \AtBeginEnvironment{lnum}{\advance\MyListDepth by 1}
    \AfterEndEnvironment{lnum}{\advance\MyListDepth by -1}
    \AtBeginEnvironment{itemize}{\advance\MyListDepth by 1}
    \AfterEndEnvironment{itemize}{\advance\MyListDepth by -1}
    \AtBeginEnvironment{enumerate}{\advance\MyListDepth by 1}
    \AfterEndEnvironment{enumerate}{\advance\MyListDepth by -1}
    % Restauración diferida: hazlo al INICIO del próximo párrafo real
    \newif\if@pendingRestore
    \@pendingRestorefalse
    \newcommand{\RequestRestore}{\global\@pendingRestoretrue}
    \everypar{%
      \if@pendingRestore
        \global\@pendingRestorefalse
        \ifInFootnote\else
          \ifnum\MyListDepth=0
            \setlength{\parskip}{\GlobalParskip}%
            \setstretch{\GlobalBaseLineStretch}%
          \fi
        \fi
      \fi
    }
    % Al final de bloques:
    \AfterEndEnvironment{equation}{\RequestRestore}
    \AfterEndEnvironment{equation*}{\RequestRestore}
    \AfterEndEnvironment{la}{\RequestRestore}
    \AfterEndEnvironment{lnum}{\RequestRestore}
    % Al final de macros:
    \apptocmd{\eqblock}{\RequestRestore}{}{}
    \apptocmd{\imagenfit}{\RequestRestore}{}{}
    
\makeatother

% ===================== ToC propio con 4 niveles (único bloque) =====================
\makeatletter

% --- Escritores: añadimos una línea al .stc en cada cabecera numerada ---
% Guardamos las versiones actuales para llamar al comportamiento normal
\let\MyTOC@orig@section\section
\let\MyTOC@orig@subsection\subsection
\let\MyTOC@orig@subsubsection\subsubsection
\let\MyTOC@orig@paragraph\paragraph

% SECTION
\renewcommand{\section}{%
  \@ifstar{\MyTOC@section@star}{\@dblarg\MyTOC@section@nostar}%
}
\def\MyTOC@section@star#1{%
  \MyTOC@orig@section*{#1}% sin entrada al .stc
}
\def\MyTOC@section@nostar[#1]#2{%
  \MyTOC@orig@section[{#1}]{#2}% comportamiento normal (escribe al .toc)
  % Escribimos también nuestra línea al .stc (texto con \numberline)
  \addcontentsline{stc}{section}{\protect\numberline{\thesection}#2}%
}

% SUBSECTION
\renewcommand{\subsection}{%
  \@ifstar{\MyTOC@subsection@star}{\@dblarg\MyTOC@subsection@nostar}%
}
\def\MyTOC@subsection@star#1{%
  \MyTOC@orig@subsection*{#1}%
}
\def\MyTOC@subsection@nostar[#1]#2{%
  \MyTOC@orig@subsection[{#1}]{#2}%
  \addcontentsline{stc}{subsection}{\protect\numberline{\thesubsection}#2}%
}

% SUBSUBSECTION
\renewcommand{\subsubsection}{%
  \@ifstar{\MyTOC@subsubsection@star}{\@dblarg\MyTOC@subsubsection@nostar}%
}
\def\MyTOC@subsubsection@star#1{%
  \MyTOC@orig@subsubsection*{#1}%
}
\def\MyTOC@subsubsection@nostar[#1]#2{%
  \MyTOC@orig@subsubsection[{#1}]{#2}%
  \addcontentsline{stc}{subsubsection}{\protect\numberline{\thesubsubsection}#2}%
}

% PARAGRAPH
\renewcommand{\paragraph}{%
  \@ifstar{\MyTOC@paragraph@star}{\@dblarg\MyTOC@paragraph@nostar}%
}
\def\MyTOC@paragraph@star#1{%
  \MyTOC@orig@paragraph*{#1}%
}
\def\MyTOC@paragraph@nostar[#1]#2{%
  \MyTOC@orig@paragraph[{#1}]{#2}%
  % Si no numeras los \paragraph, cambia \theparagraph por texto plano sin \numberline
  \addcontentsline{stc}{paragraph}{\protect\numberline{\theparagraph}#2}%
}

% --- Impresor del índice propio (lee el .stc) ---
\newcommand{\MyTOC}{%
  \par\bigskip
  {\centering\bfseries\Large Índice de contenido\par}%
  \medskip
  \begingroup
    \parindent=0pt\parskip=2pt
    \@starttoc{stc}%
  \endgroup
}

\makeatother
% ===================== fin ToC propio =====================


%%%%%%%%%%%%%%


% --- Datos del tema ---
\title{Unión Europea (VI): Política Comercial de la UE\\
\large }
\author{Marcelo Ortega}
\date{Marzo 2026}
\newcommand{\TemaCode}{3B41}

\oldtitle{
    B.43. La política comercial de la Unión Europea
}
\changerationale{
    Sin cambios.
}

\begin{document}


\maketitle
\changebox

\newpage

% --- Página de resumen y modificaciones ---
\puntosclave{ %% Escribir puntos clave Apuntes CECO
     Tradicionalmente en este tema se desarrolla en mayor medida la parte de la política comercial autónoma, aunque en la parte relativa a la política comercial convencional no se puede dejar de explicar la arquitectura de los acuerdoscomerciales modernos que firma la UE en la actualidad.
    }
\vspace{1cm}

\modificaciones{
    \textbf{NOTA (04/10/2026):} Revisar: https://x.com/judith\_arnal/status/2012803392326750313?s=20
}

\newpage
\MyTOC
\newpage

\mylistofequations
\clearpage

\MyListOfFigures
\clearpage


\pagenumbering{arabic}
\setcounter{page}{1}


\section{Introducción}



\subsection{Contextualización}


\subsubsection{Relevancia}




\subsubsection{Problemática}



\subsection{Estructura de la exposición}








\clearpage
\section{Política Comercial de la Unión: Marco jurídico-conceptual}
\puntosclave{
    Aunque la teoría económica justifica el libre comercio, la realidad es bien distinta y los países aplican políticas proteccionistas como es el caso de la UE y su PCC. 
    
    Tres motivos justifican la existencia de la PCC: (i) la integración económica conduce a un mayor bienestar económico; (ii) es una forma de avanzar en la apertura comercial; y, (iii) favorece la homogeneización de las políticas nacionales y el aumento de poder negociador.
    
    Tres características para entender el funcionamiento de la PCC: (i) es una política comunitaria; (ii) desarrollada en una unión aduanera; y, (iii) en un marco multilateral con normas comunes.
    
    La PCC es competencia exclusiva de la UE, y las decisiones se.toman por el procedimiento legislativo ordinario entre el Consejo y el Parlamento Europeo. Sin embargo, hay que considerar a los Estados miembros. 
    
    El procedimiento legislativo ordinario supone que la aprobación de los acuerdos comerciales que negocia la Comisión requiere de la aprobación conjunta del Consejo y del Parlamento Europeo.
    
    La PCC siempre debe cumplir con las normas OMC.
}


Los Estados miembros han transferido a las Instituciones de la Unión el ejercicio de buena parte de sus competencias económicas en la esfera interna en virtud de las normas referentes al mercado interior y políticas comunes y a la unión económica y monetaria. Sin embargo, las competencias asumidas por las Instituciones de la Unión para la regulación de las relaciones económicas de la UE con los países terceros son más limitadas, ya que el TFUE no contiene normas generales sobre las relaciones económicas exteriores de la UE, más allá del reconocimiento del carácter exclusivo de esta competencia de la Unión [art. 3.1.e) TFUE].

\subsection{Ámbito de aplicación}
El TFUE reglamenta parcialmente la actividad exterior de la UE en materia económica mediante los artículos $206$ a $208$, que instauran la denominada política comercial común (PCC). No obstante, las normas contenidas en estos preceptos articulan de manera deficiente la PCC y el TJUE ha tenido que mitigar en parte esta laguna, delimitando en su jurisprudencia el ámbito de aplicación y la naturaleza de las competencias atribuidas a las Instituciones de la Unión en virtud de las normas sobre PCC. A pesar de ello, la UE ha desarrollado progresivamente una PCC de carácter liberal y abierto mediante la utilización de instrumentos autónomos de política comercial y la conclusión de tratados internacionales con países terceros. Además, el Tratado de Lisboa ha supuesto un giro copernicano (EECKHOUT, 2011: 57) al permitir la plena participación del PE en este ámbito de un modo asimilado al procedimiento legislativo ordinario.

El TFUE no define la PCC ni delimita su ámbito de aplicación, ya que se limita a enumerar de forma no exhaustiva algunos instrumentos clásicos del comercio internacional que forman parte de la PCC. Inicialmente, esta enumeración de instrumentos de política comercial, unida a la vinculación establecida por el artículo $206$ entre unión aduanera y política comercial, inducía a pensar que las competencias de la Unión en esta materia se limitaban a los mecanismos clásicos del comercio internacional de mercancías.

Sin embargo, ya con el Tratado de Ámsterdam se previó la posibilidad, nunca utilizada, de ampliar por unanimidad la PCC a «negociaciones y acuerdos internacionales sobre servicios y propiedad intelectual que no estén cubiertos» por la PCC. Dando un paso más, con el Tratado de Niza se incluyó finalmente estos ámbitos en la PCC, sometidos a una decisión por mayoría cualificada y con importantes excepciones. Tras el Tratado de Lisboa, la mayoría de estas excepciones se mantienen, y así el Consejo habrá de decidir por unanimidad (art. 207.4 TFUE): cuando se trate de acuerdos en los ámbitos del comercio de servicios, o de la propiedad intelectual e industrial o sobre las inversiones extranjeras directas en donde se contengan disposiciones en las que se requiera la unanimidad para la adopción de normas internas; cuando los acuerdos en materia de servicios culturales y audiovisuales puedan perjudicar la diversidad cultural y lingüística de la Unión; y cuando los acuerdos en materia de servicios sociales, educativos o sanitarios puedan perturbar gravemente la organización nacional de dichos servicios y perjudicar a la responsabilidad de los Estados miembros en la prestación de los mismos. 

Es de enorme relevancia práctica que el Tratado de Lisboa haya incorporado las inversiones extranjeras directas al ámbito de la PCC. Esta ampliación de la PCC y la consiguiente competencia exclusiva de la UE en materia de inversiones no eliminará sin embargo la competencia estatal en este ámbito, ya que el importante número de Tratados de Protección Recíproca de Inversiones que tienen celebrados los Estados miembros con terceros regulan aspectos que quedarían fuera de la PCC. Por otro lado, la situación actual ya ha planteado la necesidad de celebrar acuerdos transitorios con esos terceros que permitan a la UE ir asumiendo esta nueva competencia.

Finalmente, el TJUE se ha visto obligado a delimitar el ámbito de aplicación de la PCC respecto de la política exterior y de seguridad, de modo que un Estado miembro unilateralmente no puede adoptar medidas sancionadoras restrictivas que puedan afectar a la PCC (asunto Centro-Com) ni la PCC puede usarse para adoptar cualquier medida restrictiva con el objeto de aplicar una posición común de la PESC (asunto Kadi).


\subsubsection{Dos posturas enfrentadas: Una delimitación jurisprudencial}
No obstante, en la práctica institucional se han planteado numerosas divergencias sobre el alcance del artículo 207 entre la Comisión y el Consejo. El TJUE ha contribuido con su jurisprudencia a la clarificación, al menos parcial, del ámbito de aplicación de la PCC, delimitando los sectores del comercio internacional cubiertos por la PCC, así como las medidas o instrumentos de política comercial incluidos en la misma.

\paragraph{Tésis finalista: Consejo}
En efecto, el Consejo ha defendido la tesis finalista, según la cual el artículo $207$ constituye la base jurídica apropiada sólo para las medidas destinadas a regular directamente los intercambios con países terceros. 

El TJUE rechazó en el Dictamen 1/78 una interpretación restrictiva del hoy artículo 207, limitada a los aspectos tradicionales del comercio exterior. La interpretación amplia se deriva pues de la enumeración no limitativa recogida en el hoy artículo 207 y del hecho de que una interpretación restrictiva provocaría distorsiones en el comercio interior como consecuencia de las disparidades de las relaciones económicas de los distintos Estados miembros con los países terceros.

El TJUE ha identificado básicamente tres criterios para realizar esta interpretación amplia del concepto de PCC: (i) el funcionamiento efectivo de la unión aduanera (sentencia Massey-Ferguson de 1973); (ii) el contenido de la PCC debe ser el mismo que el de la política comercial de los Estados (Dictamen 1/75); y, (iii) la necesidad de tener en cuenta los cambios acaecidos en la sociedad internacional en relación con los instrumentos del comercio internacional (sentencia Comisión c. Consejo de 1987).

\paragraph{Tésis instrumentalista: Comisión}
Por el contrario, la Comisión mantuvo una tesis instrumentalista u objetiva, en virtud de la cual el artículo 207 debía aplicarse siempre que se utilizara un instrumento de política comercial con independencia de la finalidad perseguida. 

Sobre la base de estos criterios, el Tribunal de Justicia ha tenido que determinar en bastantes casos si ciertos sectores del comercio internacional entraban o no en el ámbito de aplicación del artículo 207. Especial interés tiene, en este sentido, el Dictamen 1/94 en el que el TJUE se pronunció sobre la competencia de las Instituciones de la Unión para concluir los diversos acuerdos comerciales multilaterales, adoptados en la Ronda Uruguay y que figuran como anexos en el Acuerdo por el que se establece la OMC.

En este dictamen el TJUE consideró que entraban en el ámbito de la PCC todos los tratados relativos al comercio internacional de mercancías. En lo que respecta al Acuerdo General sobre el Comercio de Servicios (GATS), el Tribunal sólo admitió la inclusión en el concepto de PCC de una de las cuatro modalidades de suministro de servicios, el suministro transfronterizo, por su analogía con el comercio de mercancías. De la misma manera, en el Dictamen 1/94 el TJUE consideró, también, que los derechos de propiedad intelectual relacionados con el comercio no estaban incluidos en el ámbito de aplicación del hoy artículo 207, con alguna excepción, ya que afectan tanto o más al comercio interno que al comercio internacional. Por esta misma razón, en el Dictamen 2/92 el Tribunal entendió que no entraba dentro del ámbito de la PCC la Tercera Decisión revisada del Consejo de la OCDE, que incide en el mercado interior más que en los intercambios entre los Estados miembros y países terceros. También excluye el TJUE del artículo 207 todas las medidas referentes al comercio internacional en el ámbito de los transportes, por la existencia de una política común prevista por el TFUE en este sector.

De este modo, con la jurisprudencia derivada del Dictamen 1/94 se excluían de la PCC nuevos sectores del comercio internacional, y suponía, en cierta medida, un límite a la jurisprudencia extensiva sobre el ámbito de aplicación de la PCC. 

Con posterioridad, la Comisión ha mitigado su posición, señalando que el artículo 207 se aplica también a las medidas que, sin utilizar un instrumento de política comercial, persiguen como principal objetivo incidir en el comercio internacional.

\subsubsection{Resultados: }



\subsection{Régimen competencial: Naturaleza, alcance y contenido}
Aunque el TUE establece en la actualidad el carácter exclusivo de la PCC, la delimitación vertical de competencias en el seno del ordenamiento jurídico de la Unión ha correspondido originariamente a la jurisprudencia del TJUE. La determinación del carácter de las competencias asumidas por las Instituciones de la Unión en esta materia ha tenido una especial relevancia. En efecto, la exclusividad de la competencia de la UE elimina toda intervención estatal en la política comercial y conlleva una presión por parte de los Estados miembros para que el ámbito de aplicación de la PCC se reduzca en la mayor medida posible y para que la actuación de la Unión en las relaciones económicas con países terceros se funde en otros preceptos del TFUE, que establecen competencias compartidas de la Unión y de los Estados miembros.

\subsubsection{Naturaleza}
En lo que respecta a la naturaleza de las competencias de la Unión en materia de política comercial, el artículo 207 indica que la política comercial se basará en principios uniformes y el TJUE en su Dictamen 1/75 consagró el carácter exclusivo de la competencia de la Unión en esta materia, diciendo que la competencia paralela de los Estados quedaba excluida, porque ello supondría admitir que los Estados miembros pudieran adoptar en sus relaciones con países terceros posiciones distintas de las de la UE, que falsearían el juego institucional, quebrarían las relaciones de confianza intracomunitarias y le impedirían cumplir su función de defensa del interés común. Esta exclusividad ha sido reiteradamente confirmada por la jurisprudencia comunitaria y hoy se refleja en el artículo 3.1.e) TFUE.

El hecho de que la PCC sea una competencia exclusiva de la UE donde los Estados miembros pierden la capacidad normativa en la materia implica que las decisiones se toman a nivel de la UE mediante el procedimiento legislativo ordinario (antigua codecisión), en el que el Parlamento Europeo juega un papel fundamental desde el Tratado de Lisboa.\footnote{%
    Hasta la entrada en vigor del Tratado de Lisboa, el reparto de competencias en el ámb ito de la política comercial era algo di fuso. Aunque el Tratado de Maastricht (1993) otorgaba la competencia de la PCC a la UE. el antiguo articulo 133 del TCE fija ba unos principios uni fom1es apl icables a una enumeraGión de áreas relati vas a la política comercial. quedando sin especificar otras áreas como el comercio de servicios o las cuestiones de propiedad intelectual relacionadas con el comercio. Esta situación generó un conflicto entre la Comisión y el Consejo. ya que éste último interpretaba de manera restrictiva el texto del articulo 133, considerando fuera de la competencia exclusiva de la UE el comerc io de servicios y la propiedad intelectual. Con el Tratado de Lisboa se acaba la disputa, dando a la Comisión la competencia para todas esas áreas.
} Sin embargo, cuando los acuerdos comerciales que negocia la UE incorporan cuestiones comerciales que son competencia compartida o exclusiva de los Estados miembros, la celebración final del acuerdo comercial requiere del consentimiento de éstos, y no solo de la UE.\footnote{%
    Esto quedó con finnado tras la sentencia del Tribunal de Justicia de la UE (TJUE), del 16 de mayo de 2017, donde se dictaminó que la UE no podía celebrar el Acuerdo de Singapur por si sola deb ido a que detem1inadas partes de dicho acuerdo son objeto de una competencia compartida entre la UE y los Estados miembros. Según el TJUE, señaló que existen dos elementos del acuerdo donde la UE no tiene competencia exclusiva y necesita de las con-espondientes ratificaciones nacionales. Estos son, los tribunales de arbitraje para resolver disputas entre inversores y Estados. y las disposiciones que regulan las inversiones extranjeras de cartera.
} 


\subsubsection{Alcance}
Las dificultades encontradas en el seno del Consejo para la adopción de todas las normas necesarias para sustituir las políticas comerciales de los Estados por una PCC han exigido la modulación de la exclusividad de la competencia de la Unión en materia de política comercial, para impedir un vacío normativo que redundase en perjuicio de los intereses comerciales de la UE o de uno de sus Estados miembros. Así, dicha exclusividad se ha visto mitigada en los supuestos siguientes:
— Habilitación específica prevista en una norma de la Unión para que un Estado adopte una medida comercial (sentencia Donckerwolcke26). El TJUE también ha considerado válida una decisión del Consejo que excluía determinados productos del régimen común de exportación, permitiendo medidas comerciales autónomas de los Estados (sentencia Bulk Oil). 
— Participación de los Estados en la financiación del instrumentode política comercial. El TJUE determinó en el Dictamen 1/78 que el Acuerdo internacional sobre el caucho natural debía concluirse en forma mixta (CE y Estados miembros), porque contenía un mecanismo de intervención directa en los mercados, cuya financiación asumían los Estados miembros. Sin embargo, la aportación de fondos por parte de los Estados miembros, juntamente con la UE, para financiar el funcionamiento de una organización internacional como la OMC no impone, según el Dictamen 1/94, la conclusión mixta del tratado.

Las medidas adoptadas en estos dos casos erosionan el carácter exclusivo de la competencia de la Unión en materia de política comercial y su utilización debería ir reduciéndose para que la PCC se fuera articulando de forma definitiva y completa, como ha sucedido tras el Tratado de Lisboa con la incorporación de los ámbitos del comercio de servicios culturales y audiovisuales, así como de los servicios sociales, educativos y sanitarios.

Estos acuerdos se fundamentan, según su contenido, en el artículo 207 TFUE en solitario (materias de competencia exclusiva de la Unión) o conjuntamente en los artículos 207 y 217 TFUE (componente de asociación política, como sucede con Mercosur), y se negocian y celebran siempre conforme al procedimiento del artículo 218 TFUE. 

El tratamiento procedimental de estos acuerdos ha experimentado, sin embargo, un cambio estructural relevante a raíz del ya citado Dictamen 2/15 del TJUE, de 16 de mayo de 2017: al confirmarse que la inversión de cartera (no directa) y el mecanismo de solución de controversias inversor-Estado son materias de competencia compartida, la Comisión adoptó como práctica general, a partir de ese momento, la de escindir en dos instrumentos jurídicos separados lo que antes se negociaba como un único acuerdo mixto: de un lado, un acuerdo comercial en sentido estricto, de competencia exclusiva de la Unión, que puede entrar plenamente en vigor tan pronto como lo aprueban el Consejo y el Parlamento Europeo, sin necesidad de ratificación por los Estados miembros; de otro, un acuerdo de protección de inversiones, de competencia compartida, que sí exige esa ratificación nacional y que, en consecuencia, tarda años --a veces indefinidamente-- en entrar en vigor. Los acuerdos con Corea del Sur y con Canadá, anteriores al Dictamen, siguen el modelo antiguo de acuerdo único mixto; los acuerdos con Singapur y con Vietnam son los primeros en adoptar el modelo escindido, que se ha consolidado desde entonces como estándar.

Este acuerdo protagonizó, además, el primer caso de activación efectiva del mecanismo de solución de controversias del capítulo TSD en la práctica convencional de la Unión. A raíz de una denuncia presentada por la Confederación Europea de Sindicatos y otras organizaciones sindicales, un panel de expertos independiente concluyó, el 25 de enero de 2021, que Corea del Sur había incumplido su compromiso, asumido en el propio acuerdo, de desplegar esfuerzos continuos y sostenidos para ratificar los convenios fundamentales de la Organización Internacional del Trabajo.\footnote{%
    Panel de expertos constituido en el marco del Acuerdo de Libre Comercio UE-Corea del Sur, dictamen de 25 de enero de 2021.
} Aunque el pronunciamiento del panel no llevaba aparejada la posibilidad de sanciones comerciales, la presión política derivada de esta decisión contribuyó a que Corea del Sur ratificara, a lo largo de 2021, los convenios de la OIT sobre libertad sindical y negociación colectiva (C87 y C98) y sobre trabajo forzoso (C29). Este episodio se cita habitualmente en la doctrina como la prueba de que los capítulos TSD, pese a su carácter originariamente no sancionador, pueden generar un efecto de presión reputacional con consecuencias prácticas tangibles.


\subsubsection{Contenido: Los intrumentos de Política Comercial Común}
La larga andadura de la Unión como actor económico internacional y su activa participación en la regulación de las relaciones económicas internacionales ha desembocado en que la UE disponga de un amplísimo instrumental motivado no sólo por su consolidada experiencia histórica, sino también por el extraordinario desarrollo que han experimentado estas últimas en épocas recientes. Las Instituciones de la Unión han ido elaborando sobre la base de los artículos 206 y 207 TFUE una política comercial común que cuenta con un amplio arsenal de instrumentos regulados en los numerosos tratados multilaterales y bilaterales de la UE con países terceros y, también, en normas de Derecho derivado de la Unión. 

En el estudio de los instrumentos de la PCC era clásica la distinción entre las llamadas PCC autónoma —que agrupaba aquellas medidas unilaterales adoptadas por la UE— y la PCC convencional —término con el que se designaban los tratados comerciales concluidos con países terceros y destinados a regular sus relaciones comerciales—. Sin embargo, no parece posible mantener la distinción entre instrumentos convencionales y autónomos de
política comercial, porque en muchos casos la interacción entre ambos es muy directa y, fundamentalmente, porque los instrumentos autónomos tienen un contenido predeterminado por tratados internacionales —especialmente, del sistema OMC—. 

: cambios arancelarios, celebración de acuerdos arancelarios y comerciales relativos al comercio de mercancías y de servicios, los aspectos comerciales de la propiedad intelectual, la inversión extranjera directa, la uniformización de las medidas de liberalización, la política de exportación y las medidas de defensa comercial, entre ellas las adoptadas en caso de dumping o subvenciones.  No obstante, el propio artículo $207.4$ TFUE introduce una excepción relevante: se exige unanimidad para la negociación y celebración de acuerdos en el ámbito del comercio de servicios culturales y audiovisuales, cuando exista riesgo de perjuicio para la diversidad cultural y lingüística de la Unión, así como en el ámbito del comercio de servicios sociales, de educación y de salud, cuando exista riesgo de perturbación grave de la organización de dichos servicios a nivel nacional y de menoscabo de la responsabilidad de los Estados miembros para prestarlos. Esta excepción explica, por ejemplo, la exclusión sistemática de ciertos servicios audiovisuales de los capítulos de servicios en los ALC de Nueva Generación.



En realidad, \textbf{los instrumentos de política comercial admiten tres modalidades de aplicación: autónoma, convencional y preferencial} (GONZÁLEZ ALONSO, 2000: 428). El régimen convencional es el más importante, ya que comprende las normas aplicables a los países miembros de la OMC, que son la gran mayoría y que concentran la práctica totalidad del comercio internacional. El régimen autónomo de la política comercial, aplicado a los países de comercio de Estado es hoy muy residual, máxime desde la incorporación de la República Popular China a la OMC. El régimen preferencial agrupa a aquellos Estados a los que la UE otorga preferencias comerciales respecto a las acordadas en el seno de la OMC. Los regímenes preferenciales de la Unión son muy numerosos y con distinto fundamento jurídico. Según el nivel de concesión arancelaria, irían desde el sistema UE de preferencias arancelarias generalizadas, hasta el Acuerdo EEE de 1992, por el que se extiende el mercado interior a los países de la EFTA (Noruega, Islandia y Liechtenstein). 

Para la exposición de cuáles son los instrumentos más importantes con los que cuenta la política comercial común en la actualidad, seguiremos la sistematización de LÓPEZ ESCUDERO (2003: 59-63) que distingue entre aplicables a la importación, a la exportación, instrumentos de defensa comercial y regímenes comerciales preferenciales:


















La PCC posee tres características que penniten comprenqer su funcionamiento: i) es una política comunitaria con unos objetivos bien definidos; ii) que se desarrolla en una unión aduanera; y iii) en un marco multilateral con nonnas comunes.

Los objetivos de la PCC se establecen en el TFUE. En su artículo 206 se menciona que «mediante el establecimiento de una unión aduanera, la UE contribuirá, en el interés común, al desarrollo armonioso del comercio mundial, a la supresión progresiva de las restricciones a los intercambios internacionales y a las inversiones extranjeras directas, así como a la reducción de las barreras arancelarias y de otro tipo». Adicionalmente, la PCC es fundamental desde otrospuntos de vista. constituyéndose como uno de los pilares de la acción exterior de la UE (artículo 207 del TFUE) y como un instrumento coadyuvante en otros ámbitos.\footnote{%
    Es el caso de la política de cooperación al desarrollo, que aplica instrumentos comerciales. como el SPG o el Aidfor Tra de. para lograr sus objetivos: la política agrícola común (PAC), que se basa en el principio de preferencia comun itaria para dar una mayor protección a los productos agroalimentarios de la UE frente al exterior; o la política de competenci a. donde las medidas compensatorias o derechos antidumping pueden apoyar sus fines.
} 

En el caso de la negociación de nuevos acuerdos comerciales, la cadena de decisión es la siguiente (artículos 207 y 218 del TFUE): en primer lugar, la Comisión Europea (en adelante, Comisión) presenta al Consejo sus recomendaciones para iniciar las negociaciones que considere necesarias. A continuación, el Consejo autoriza la apertura de las negociaciones. y da un mandato negociador de carácter vinculante a la Comisión, en el que se fijan las directrices para la negociación. Las negociaciones son llevadas por la Comisión, en consulta con el Comité de Política Comercial, al que deberá infonnar periódicamente, y el cual la asistirá en la materia. También se deberá informar al Parlamento Europeo de la marcha de las negociaciones. 

Al finalizar las negociaciones, la Comisión presenta el texto al Consejo de Asuntos Exteriores-Comercio, que debe votar para aceptarlo o rechazarlo.\footnote{%
    En este caso. el Consejo de Asuntos Exteriores-Comercio lo preside el ministro de tumo del Estado miembro que preside el Consejo, al tratarse de un tema comercial, y no el Alto Representante.
} Al ser una competencia exclusiva de la UE, la regla de votación que se aplica en el seno del Consejo es la de la mayoría cualificada (doble mayoría, hoy vigente) si se trata de comercio de mercancías, comercio de servicios, aspectos de los derechos de propiedad intelectual relacionados con el comercio e IED. Sin embargo, existen excepciones para las cuales el Consejo debe pronunciarse por unanimidad: el comercio de servicios culturales y audiovisuales, cuando los acuerdos puedan perju dicar a la diversidad cultural y lingüística de la UE; y el comercio de servicios sociales, educativos y sanitarios, cuando los acuerdos puedan afectar gravemente a la organización nacional de dichos servicios y perjudicar a la responsabilidad de los Estados miembros en la prestación de los mismos. Finalmente, en caso de un voto positivo por el Consejo para la finna del acuerdo, el Parlamento Europeo debe aprobar el texto también. Una vez recibido el visto bueno del Parlamento Europeo, la Comisión puede firmar el acuerdo en nombre de la UE. 

Fuera de la competencia exclusiva de la UE queda la política de promoción exterior, referida al apoyo genérico a la internacionalización de la empresa mediante diferentes tipos de instrumentos, ya sean informativos, de formación, o de financiación (apoyo financiero, misiones comerciales, etc.). No obstante, a pesar de ser una competencia de índole nacional, existe una exigencia de armonización progresiva de los instrumentos utilizados por los Estados miembros en este ámbito, a fin de evitar que se falsee la competencia entre empresas de fa UE en terceros mercados.


\clearpage
\section{Política Comercial Común: Importaciones y Exportaciones}




\subsection{Política Comercial Común: Elementos comunes}


\subsubsection{Régimen Comercial General: Principio de libertad comercial}
El régimen de libertad, o Régimen General, constituye la regla base del sistema y su punto de partida lógico y jurídico: la importación y la exportación de mercancías no están sometidas a restricción cuantitativa alguna ni a requisito de autorización previa, de suerte que la operación se perfecciona con la sola presentación de la declaración aduanera conforme al Código Aduanero de la Unión (Reglamento (UE) 952/2013) y el pago de los derechos del arancel aduanero común y demás tributos exigibles. 

Los artículos primeros de los Reglamentos 2015/478 y 2015/479 lo proclaman expresamente al disponer que las importaciones y exportaciones serán libres y no estarán sometidas a restricciones cuantitativas, sin perjuicio de las excepciones que los propios reglamentos habilitan. 

El fundamento de este Régimen es doble. Por un lado, hacia el exterior, materializa las obligaciones asumidas por la Unión en la OMC, singularmente la prohibición de restricciones cuantitativas del art. XI del GATT y la consolidación arancelaria, que convierten el arancel en el único instrumento ordinario de protección frente a los demás miembros. Por otro lado, hacia el interior, expresa la visión liberal del constituyente europeo plasmada en el artículo $206$ TFUE. 

De hecho, un dato frecuentemente omitido revela hasta qué punto la libertad es hoy el régimen universal y no meramente el convencional: el Reglamento 2015/755, aplicable a las importaciones procedentes de los escasos países ajenos a la OMC (Bielorrusia, Corea del Norte y algún otro), consagra idéntico principio de libertad de importación. La Unión, que frente a estos países no está convencionalmente obligada a nada, liberaliza unilateralmente; el antiguo régimen autónomo aplicado a los países de comercio de Estado ha quedado así reducido a una diferencia de base normativa sin apenas diferencia de contenido material. 

\subsubsection{Regímenes especiales: Derogación de la norma}
Ahora bien, uno de los principales objetivos de todo Estado de derecho es proteger la seguridad nacional e internacional, impedir el tráfico ilícito y la proliferación de armamento y tecnologías sensibles que puedan ser utilizadas por actores que amenacen la paz y la seguridad, o que estén involucrados en actividades terroristas. En este sentido, resulta fundamental garantizar que tanto importaciones como exportaciones sean coherentes con los compromisos internacionales adquiridos por España, cumpliendo con el Derecho Internacional. Esto implica asegurar que no contribuyan a la violación de los derechos humanos, ni a la intensificación de conflictos armados, ni fomenten nuevas tensiones o situaciones de pobreza. Además, en lo que se refiere a determinados productos, resulta frecuentemente necesario estudiar las condiciones de las exportaciones e importaciones, su evolución y los distintos aspectos de la situación económica y comercial de la forma más rápida posible, incluso si esta fuese utilizado para la adopción de medidas concretas (salvaguardias por ejemplo). Para estos productos puede resultar necesario someter determinadas importaciones a una vigilancia de la Unión.

Por todo lo expuesto, el Derecho de la Unión contempla una serie de regímenes especiales que obligan a adoptar un régimen especial de vigilancia en virtud de los objetivos expuestos y siempre para productos tasados. Esto convierte el régimen comercial de la Unión en una estructura binaria: la libertad como régimen general, la vigilancia como mecanismo de seguimiento que no restringe el comercio pero lo somete a observación documental. 

Como la Unión carece de una administración propia de comercio exterior que expida licencias y certificados a los operadores, pues tampoco cuenta con una Administración de Aduanas propia, la aplicación descansa, conforme al modelo general de administración indirecta del Derecho de la Unión, sobre las autoridades nacionales. En este sentido, la respuesta administrativa Española fue reconducir toda la dispersión normativa europea a cuatro regímenes operativos según la intensidad de la intervención, siendo el más leve (ausente de intervención) el que acabamos de ver.

\paragraph{Régimen de vigilancia}
El primero de estos, por intesidad, es el régimen de vigilancia.

El régimen de vigilancia somete la importación de determinados productos (y excepcionalmente su exportación) a la presentación previa de un documento, cuya nota definitoria es la automaticidad de su expedición: la autoridad nacional competente debe emitirlo gratuitamente y en un plazo de pocos días hábiles para cualquier cantidad solicitada, sin margen alguno de apreciación y sin posibilidad de denegación. De ahí que la vigilancia no constituya materialmente una restricción al comercio, lo que la hace compatible sin dificultad con el artículo XI del GATT, ya que se trata de un instrumento de información, no de protección. 

Normalmente, la razón argumental es la insuficiencia estadística ordinaria, que llega a la Comisión con meses de retraso, cuando la evolución de las importaciones de un producto sensible puede requerir reacción en semanas. De esta forma, la información recaba mediante los documentos proporciona a la Comisión conocimiento en tiempo casi real de los volúmenes, precios y orígenes de las importaciones del producto vigilado, y cumple con ello una doble función: de alerta temprana, permitiendo detectar incrementos súbitos de importaciones, y de preconstitución probatoria, pues los datos recabados sirven para documentar el perjuicio grave o la amenaza de perjuicio grave a los productores de la Unión que constituye el presupuesto de una eventual medida de salvaguardia. La vigilancia es, en este sentido, la antesala procedimental natural de la salvaguardia. Se regula en los artículos 10 a 14 del Reglamento 2015/478 (con previsión paralela en el Reglamento 2015/755) y puede establecerse erga omnes o circunscrita a orígenes determinados, siempre respecto de productos concretos designados por la Comisión cuando la evolución del mercado lo aconseje. 

\paragraph{Régimen de certificación}
El segundo en intensidad sería el Régimen de Certificación (que ocupa una posición intermedia entre la vigilancia y la autorización).

Su nota definitoria es el carácter reglado de la expedición: la operación exige la presentación de un certificado o licencia, pero el operador que acredita el cumplimiento de los requisitos objetivos establecidos por la norma tiene derecho a obtenerlo, sin que la Administración disponga de margen de oportunidad para denegarlo. No hay, pues, ponderación discrecional de intereses, sino verificación documental de condiciones. 

El régimen cumple tres funciones distintas que conviene deslindar:
\begin{la}
    \item \textbf{Previsión}. La primera es la ordenación y previsión de flujos, cuyo territorio natural es la PAC: los certificados \textbf{AGRIM} para la importación y \textbf{AGREX} para la exportación de determinados productos agrícolas (arroz, ajos, cáñamo), regulados por el Reglamento Delegado (UE) 2016/1237 y el Reglamento de Ejecución (UE) 2016/1239, presentan la particularidad de expedirse previa constitución de una garantía que se ejecuta si la operación no se realiza durante el periodo de validez del certificado, porque el certificado no confiere solo un derecho sino que impone la obligación de importar o exportar. Por tanto, mediante esta técnica el sistema obtiene una previsión fiable de los flujos físicos que van a producirse, información esencial para la gestión de los mercados agrarios y de los contingentes arancelarios;
    \item \textbf{Trazabilidad}. La segunda función es la trazabilidad exigida por compromisos internacionales de naturaleza ética o de seguridad. Por ejemplo, sería el caso del sistema de certificación del Proceso de Kimberley para los diamantes en bruto, aplicado por el Reglamento (CE) 2368/2002 para excluir del comercio los llamados diamantes de conflicto, que solo opera respecto de los países participantes en el Proceso;
    \item \textbf{Legalidad}. La tercera es la acreditación de la legalidad de origen por razones ambientales.\footnote{%
        Por ejemplo, las licencias FLEGT del Reglamento (CE) 2173/2005, que certifican la tala legal de la madera importada de países socios y que se encuentran en trance de absorción por el sistema de diligencia debida del Reglamento de deforestación (UE) 2023/1115; los permisos CITES del Reglamento (CE) 338/97 para el comercio de especímenes de fauna y flora silvestres amenazadas; y los certificados de captura del Reglamento (CE) 1005/2008 contra la pesca ilegal, no declarada y no reglamentada.
    }
\end{la}

A diferencia de la vigilancia y la autorización, que suelen definirse por producto y origen, la certificación responde en general a la naturaleza de la mercancía y opera erga omnes (aunque algunos sí sobre países participantes o socios determinados, Kimberley y FLEGT). Su relevancia es alta y creciente, no tanto por el volumen de comercio cubierto cuanto por su valor de modelo: el esquema certificación-trazabilidad-diligencia debida es la técnica sobre la que se está construyendo la nueva generación de comercio regulado por condicionalidad ambiental y de derechos humanos, de la que el Reglamento de deforestación es la generalización más ambiciosa y que veremos en la segunda sección. Procedimentalmente, el certificado se solicita ante el organismo nacional competente (en España, según el producto, la Secretaría de Estado de Comercio o el Fondo Español de Garantía Agraria), con constitución de garantía en el caso agrícola, se expide de forma reglada en plazos breves y se presenta en aduana como documento de acompañamiento de la declaración.

\paragraph{Régimen de autorización administrativa}
El tercero y más intensivo sería el \textbf{Régimen de Autorización}.

En realidad, es el único genuinamente restrictivo de los tres regímenes específicos, pues la operación queda sometida a una autorización previa cuya concesión es discrecional, en el sentido de que la Administración pondera caso por caso circunstancias como el destino, el usuario final y el uso declarado de la mercancía, y puede denegar motivadamente la operación (aplicándose incluso silencio negativo). Como se puede apreciar, la libertad comercial cede de verdad, y cede porque concurren intereses superiores que tanto el Derecho de la Unión como el Derecho OMC reconocen como títulos habilitantes de la restricción: la seguridad internacional y la no proliferación, las obligaciones de política exterior, etc. Motivos, de hecho, reconducibles a las excepciones generales y de seguridad de los arts. XX y XXI del GATT, lo que explica que estas restricciones se apliquen indistintamente a miembros y no miembros de la OMC. 

Su ámbito principal es la exportación, y su expresión más desarrollada es el control del comercio de productos y tecnologías de doble uso civil y militar, regulado por el Reglamento (UE) 2021/821, que refundió y modernizó el régimen anterior incorporando controles sobre tecnologías de cibervigilancia: el reglamento opera sobre listas de productos controlados y articula un sistema graduado de autorizaciones individuales, globales y generales de la Unión en función del producto, el destino y la fiabilidad del exportador. Junto al doble uso se sitúan el comercio de material de defensa propiamente dicho y otras materias reguladas (bienes susceptibles de utilizarse para la pena capital o la tortura).

Con un peso creciente que ha transformado la fisonomía del régimen, las medidas restrictivas adoptadas sobre la base del artículo $215$ TFUE en ejecución de decisiones PESC, cuyo ejemplo paradigmático son los sucesivos paquetes de sanciones contra Rusia desde 2014 e intensivamente desde 2022, en los que la técnica autorizatoria convive con prohibiciones directas de exportación e importación por sectores enteros. En el ámbito de la importación, la técnica autorizatoria sirve además para administrar los contingentes derivados de medidas de salvaguardia y para controlar mercancías singulares como los precursores de drogas o los traslados de residuos.

La versatilidad y rapidez de despliegue del régimen quedaron demostradas con el mecanismo temporal de autorización de exportación de vacunas contra la COVID-19 adoptado en 2021 sobre la base del Reglamento 2015/479, que acreditó que la libertad de exportación puede excepcionarse en días cuando concurre un riesgo de desabastecimiento. En términos de relevancia, es el régimen que más ha crecido en la última década, porque es el cauce por el que discurre la agenda de seguridad económica de la Unión, que veremos más adelante. Procedimentalmente, la solicitud se presenta ante la autoridad nacional competente —en España, la Secretaría de Estado de Comercio, con intervención de la JIMDDU en doble uso y defensa—, se evalúa conforme a los criterios materiales del reglamento aplicable, y la resolución, que puede imponer condiciones o denegarse motivadamente, suele acompañarse de obligaciones documentales ulteriores como los certificados de último destino.

\subsubsection{Obligaciones comunes}


\paragraph{Registro: EORI}
Para entablar relaciones comerciales con el exterior de la Unión se debe solicitar en la administración aduanera competente el denominado número de registro e identificación de operadores económicos (EORI).

El EORI es un número de identificación válido en toda la Unión y obligatorio para el despacho de todo tipo de operaciones aduaneras en el territorio aduanero de la Unión, como la importación, la exportación y el tránsito. Los números EORI identifican de manera única a los operadores económicos y otras personas. En cualquier momento, a una persona solo se le puede asignar un número EORI válido. Los operadores económicos deben comunicar este número a las autoridades aduaneras de los países de la UE para las operaciones aduaneras.

¿Quién necesita un número EORI? De base, cualquier operador económico establecido en el territorio aduanero de la UE, pero también operador económico que no están establecidos en el territorio aduanero de la Unión, pero llevan a cabo las siguientes actividades aduaneras: presentación de declaraciones en aduana, presentación de declaraciones sumarias de entrada (ENS), presentación de declaraciones sumarias de salida (EXS), presentación de declaraciones de depósito temporal, actuando como transportista. Además, también personas distintas a los operadores económicos si así lo exige la legislación nacional del país de la UE en cuestión o otros actos jurídicos de la UE, o si participan en actividades para las que debe facilitarse un número EORI. 

\paragraph{Tributarias: Fiscalidad Indirecta}

A efectos tributarios, dentro de las importaciones y exportaciones, la materia fundamental y de enorme casuística, es el tratamiento del IVA.

Siendo empresario, cuando se importa, hay que diferenciar si se trata de un bien o un servicio. Al importar un bien, se pagará el arancel correspondiente en aduana (dependiendo del producto y el país de origen) y el IVA a través del denominado modelo 031.

Tanto en las exportaciones como importaciones, hay que cumplimentar el famoso y tantas veces citado DUA (Documento Único Administrativo) que es lo que va a permitir deducir las cuotas soportadas de IVA en nuestra declaración.


\subsection{Instrumentos aplicables a la importación: Código Aduanero Común}
La Unión Aduanera Común se estableció el 1 de julio de 1968, un año y medio año antes de lo previsto, y supuso el establecimiento de un Arancel Exterior Común, así como la entrada en vigor de un corpus normativo propio que permitiera el tránsito de mercancías entre los Estados miembros. Más tarde, con el Acta Única Europea (1986), se fija como objetivo la Unidad Funcional del control exterior de mercancias, pues el proyecto comunitario resulta singular en tanto que mantiene la gestión de Aduanas en manos de las Autoridades Nacionales. Con este gran paso, imprescindible para el desarrollo del Mercado Único Interior, se consigue alcanzar la perfecta unidad funcional de cara al exterior del Mercado Interior de la Unión.\footnote{%
    La evolución de este proceso, por su relevancia para el MIU se trata más en detalle en el [\authorfont{Ver Tema 3.B.39}]
}

De este proceso de unificación, lo que más nos interesa en tanto tratamos la PCC es su incidencia económica y su proyección comunitaria. De esta forma, de entre todos los Actos y Decisiones adoptados durante el breve proceso que transcurre desde el AUE hasta la formalización del MIU, prevalece en importancia sobre cualquier otro el \textbf{Código Aduanero Común}. Adoptado en 1992, este Cógido contiene todas las disposiciones relativas a la importación de mercancías procedentes de terceros Estados, como la configuración y aplicación del AEC y del régimen general de importación en la UE. 

Esta libertad de importación sólo puede verse limitada según el Reglamento n.º 2015/755 en supuestos de aumento de importación de un producto que pueda causar un perjuicio grave a los productores de la UE. En estos casos, el Reglamento n.º 2015/75 da poderes a la Comisión para llevar a cabo una investigación y adoptar, en su caso, medidas de vigilancia o medidas de salvaguardia de conformidad con el derecho de la OMC. Estas medidas de salvaguardia, que son siempre temporales, se adoptarán cuando haya importaciones de un producto en tales cantidades o condiciones que ocasione o pueda ocasionar un perjuicio grave a los productores de la Unión de productos similares o directamente competitivos. Las medidas de salvaguardia suelen consistir en el establecimiento de contingentes o restricciones cuantitativas limitadoras de la cantidad de mercancías que pueden importarse. 

Se trata de la última manifestación del mandato previsto en el art. 28.1 TFUE y, en la actualidad, después de varias modificaciones, el Código Aduanero de la Unión está recogido en el Reglamento n.º 952/2013, aplicable en su totalidad desde el 1 de junio de 2016.\footnote{
\textbf{NOTA.-} El territorio que comprende la Unión Aduanera excede los límtes territoriales de la Unión. Es el caso de Turquía (en 1995) o microestados europeos, como Mónaco (1968), Andorra y San Marino. No obstante, los ámbitos materiales son más restringidos. Por ejemplo, en el caso de Turquía solo abarca los productos industriales y productos agrícolas transformados, y para productos agrícolas, carbón y acero existen algunas concesiones comerciales recíprocas. En diciembre de 2016, la Comisión propuso ampliar las relaciones comerciales bilaterales con Turquía en el marco de la unión aduanera a ámbitos como los servicios, la contratación pública y el desarrollo sostenible; aunque el Consejo aún no ha adoptado el mandato. Por otro lado, el Acuerdo Andorra-UE excluye también los productos agrícolas. No obstante, también sucede al revés. Hay territorios que no forman parte de la Unión Aduanera. Por ejemplo, Ceuta y Melilla.

El estatuto más peculiar lo ostenta la Ciudad del Vaticano: no forma parte de la Unión Aduanera, pero se encuentra ligado a Italia a través de los \href{https://es.wikipedia.org/wiki/Pactos_de_Letrán}{Pactos de Letrán (1929)}, que que extiende al Vaticano las condiciones de la Unión Económica y Monetaria.
} Véamos sus componentes más en detalle.

\subsubsection{Arancel Exterior Común: TARIC y NC}
Los derehos de aduana de la Unión, se establecen de conformidad con los tipos impositivos recogidos en el Arancel Aduanero Común (AAC) y se configura como el principal instrumento sobre las importaciones de la PCC de la Unión. Con una larga historia, que se remonta al uno de julio de 1968, constituye el primer instrumento de la PCC.\footnote{%
    Además de suponer una importante fuente de ingresos para el Marco Financiero Plurianual [\authorfont{Ver Tema 3.B.37}]. Los Estados miembros retienen en concepto de gastos de recaudación el $25\%$ del valor (lo que se conoce como el efecto Rotterdam), sirviendo incentivo para garantizar una recaudación diligente de los importes adeudados.
}

El AAC se calcula para cada tipo de mercancías, en función de su origen y de su valor.\footnote{%
    Incialmente, para ponerlo en marcha, se calculó como la media de los tipos vigentes de los aranceles nacionales, y que debían adoptar los países que se iban adhiriendo a la Unión, con un breve periodo de transición. Se modificó ligeramente en el Reglamento 2658/87 con el fin de unificar las funciones estadísticas y de nomenclatura. Este Reglamento estableció la base jurídica para la posterior configuración del Sistema TARIC, introduciendo el sistema común de codificación y clasificación de mercancías, la Nomenclatura Combinada (NC), esencial para el tratamiento de los datos y la publicación de estadísticas comerciales de la UE.
} En esta labor resulta esencial tanto el \href{https://taxation-customs.ec.europa.eu/customs/common-customs-tariff-cct/tariff-classification-goods/eu-customs-tariff-taric_en}{\textbf{Sistema TARIC}}, o arancel integrado de la Unión Europea, que es una base de datos que integra todas las medidas relacionadas con el AAC y la legislación comercial y agrícola; como la \href{https://taxation-customs.ec.europa.eu/customs/common-customs-tariff-cct/tariff-classification-goods/combined-nomenclature_en}{\textbf{Nomenclatura Combinada}} (NC), resultado de la fusión de las entonces existentes nomenclaturas del arancel aduanero común y la NIMEXE (Nomenclatura estadística de la UE) y creada para cumplir los requisitos del AAC y las estadísticas comerciales de la UE, que actualmente actúa de herramienta de clasificación de mercancías.\footnote{%
    Se trata de un nuevo desarrollo de la nomenclatura del Sistema Armonizado de la Organización Mundial de Aduanas. La NC se utiliza para clasificar la mayoría de las mercancías cuando se declaran en aduana en la UE. Las subpartidas NC indicadas en las declaraciones de mercancías importadas y exportadas determinan: qué tipo de derecho de aduana se aplica y cómo se tratan los bienes con fines estadísticos o para otras políticas de la UE
}

El nivel de los derechos de aduana es bastante bajo, como consecuencia de las negociaciones desarrolladas en el seno de la OMC y de las rebajas arancelarias acordadas adicionalmente por la UE a países terceros en su acuerdos comerciales preferenciales. En la actualidad, el tipo arancelario medio para mercancías no agrícolas se encuentra en tomo al $3,5\%$, mientras que subsisten notables crestas arancelarias para productos agrícolas, siendo el tipo arancelario medio para éstos de alrededor del $13\%$.

\paragraph{Circulación intracomunitaria (\textit{ex} art. 3 TFUE): Despacho a libre práctica}
El \textbf{despacho a libre práctica}, que permite a las mercancías circular libremente por la Unión como si fueran mercancías originarias de la Unión y se establece como principio general la no aplicación a las importaciones por parte de las autoridades de la UE y de sus Estados miembros de ningún tipo de restricción, excepción hecha de los derechos de aduana. 

\subsubsection{Excepciones: Exenciones circunstanciales del AAC}
Existen una serie de excepciones que llevan a que el AAC se reduzca o no se aplique, atendiendo a dos circunstancias.

\paragraph{Medidas de Economía Arancelaria}
La primera de ellas es lo que se conoce como \textbf{medidas de economía arancelaria}. Se trata de un conjunto de instrumentos que emplea la Unión para modificar el AAC con el objetivo de corregir desequilibrios coyunturales (como una producción interior insuficiente de alguna mercancía) o favorecer determinadas importaciones por otras razones económicas, culturales o científicas. La heterogeneidad de razones que conducen a la aplicación de este régimen especial, se manifiesta también en los medios materiales, donde se contemplan varias modalidades:
\begin{la}
    \item \textbf{Contingentes arancelarios}. Los contingentes arancelarios permiten la importación de una cantidad determinada de una mercancía extracomunitaria, por un periodo específico de tiempo, con un derecho arancelario inferior o nulo respecto al establecido con carácter general. Las condiciones bajo las cuales se puede aplicar un contingente están tasadas y, normalmente, afectarán a la importación de materias primas, productos semielaborados o componentes que no están disponibles en la UE, en absoluto o en cantidad suficiente. No se concederá ningún contingente arancelario para los productos terminados;
    \item \textbf{Suspensión arancelaria}. La suspensión arancelaria consiste en dejar de aplicar, total o parcialmente, los derechos arancelarios previstos, siempre de forma temporal. A diferencia de los contingentes, no existe un límite de cantidad de mercancía que puede beneficiarse de esta exención y, igual que en el caso anterior, deben cumplirse una serie de condiciones. En algunos casos, la concesión de la suspensión está condicionada al cumplimiento de un determinado uso o destino de la mercancía; y,
    \item \textbf{Franquicias arancelarias}. La franquicia arancelaria es una exención total de derechos aduaneros que se diferencian de los contingentes y suspensiones porque no tienen ninguna limitación temporal. Como ejemplo suele citarse las mercancías contenidas en el equipaje personal de los viajeros procedentes de un país no comunitario, siempre que no tengan un fin comercial.
\end{la}

\paragraph{Regímenes Aduaneros Especiales}
Otra modalidad particular serían los regímenes económicos aduaneros que persiguen matener la competitividad de los productos comunitarios. Previstos en el Código Aduanero Común, se trata de una exención arancelaria sobre bienes que están sujetos a reimportación/reexportación, desde la Unión a un país tercero, con el objetivo de llevar a cabo un proceso parcial de transformación. En este ámbito la doctrina distingue entre cuatro modalidad.

La primera de ellas es el \textbf{depósito y almacenamiento de mercancías}. Se trata de dos regímenes aduaneros que permiten el almacenamiento de mercancías no pertenecientes a la Unión, sin que estén sujetas a derechos de importación, otros gravámenes prescritos por otras disposiciones pertinentes en vigor y a medidas de política comercial en la medida en que no prohíban la entrada de mercancías en el territorio aduanero de la Unión o su salida de él. Los regímenes especiales de depósito se dividen en dos: el régimen especial de \textbf{depósito aduanero} y el régimen especial de \textbf{zona franca}.\footnote{%
    La expresión depósito aduanero tiene una doble acepción: (i) como instalación o inmueble, el depósito aduanero es todo lugar reconocido por las autoridades aduaneras y sometido a su control, en el que puedan almacenarse mercancías de importación en las condiciones establecidas en la normativa aduanera; o, (ii) como una situación jurídico-aduanera en la que se encuentran las mercancías, el régimen especial de depósito aduanero es aquel que permite el almacenamiento, en unas instalaciones de depósito aduanero, de mercancías no pertenecientes a la Unión en las condiciones establecidas para los regímenes especiales de depósito. Normalmente ambas acepciones suelen utilizarse de manera simultánea, ya que se solicita que sea autorizado un local como depósito aduanero, para así poder vincular al régimen especial de depósito aduanero mercancías procedentes de terceros países en dichas instalaciones.
} En el depósito aduanero sólo se contempla la posibilidad de almacenar, en las instalaciones autorizadas, mercancías no pertenecientes a la Unión. No obstante, las autoridades aduaneras pueden incluir, en el mismo acuerdo, una autorización para almacenar conjuntamente mercancías vinculadas al régimen de depósito aduanero con mercancías de la Unión, siempre que exista una necesidad económica y no quede comprometida la vigilancia aduanera. Por otro lado, las zonas francas, que no aparecen definidas expresamente en el Código Aduanero de la Unión, se podrían definir como partes cerradas del territorio aduanero de la Unión, separadas del resto del mismo, en las que se pueden almacenar toda clase de mercancías no pertenecientes a la Unión, sin estar sujetas a derechos de importación, gravámenes prescritos por otras disposiciones pertinentes en vigor y/o medidas de política comercial en la medida en que no prohíban la entrada de mercancías en el territorio aduanero de la Unión o su salida de él. Son los Estados miembros quienes designan, según su voluntad, que determinadas partes del territorio aduanero de la Unión tengan la consideración de zonas francas. Los Estados miembros deben comunicar a la Comisión Europea información sobre las zonas francas que estén en funcionamiento en sus respectivos territorios, que deberán estar cercadas. El perímetro y los puntos de acceso y de salida de las zonas francas estarán sometidos a vigilancia aduanera, pudiendo ser sometidos a controles aduaneros, tanto las mercancías como los medios de transporte que entren o salgan de la zona franca.

La segunda de éstos sería el régimen especial por \textbf{destinos especiales}. En el marco del régimen de importación temporal, las mercancías no pertenecientes a la Unión destinadas a la reexportación podrán ser objeto de un destino especial en el territorio aduanero de la Unión con exención total o parcial de derechos de importación y sin estar sometidas a lo siguiente a otros gravámenes prescritos por otras disposiciones pertinentes en vigor o medidas de política comercial en la medida en que no prohíban la entrada de mercancías en el territorio aduanero de la Unión o su salida de él, sin que esté previsto que las mercancías sufran cambio alguno. En estos casos existen dos formas de sujetarse al régimen: \textbf{importación temporal} o \textbf{destino final}. Mientras la importación temporal es muy sencilla, puesto que el ejemplo clásico es cualquier importación que tenga lugar con motivo de un evento ferial, el destino final es un poco más complejo: las mercancías pueden despacharse a libre práctica, con exención de derechos o con un tipo reducido cuando van a ser utilizadas para un uso o destino específico. Por tanto, este último régimen lleva consigo una vigilancia aduanera posterior a la importación, que finalizará cuando las mercancías se hayan destinado a los fines establecidos para la aplicación de la exención de derechos o de tipo reducido, hayan salido del territorio aduanero de la Unión, destruidas o abandonadas en favor del Estado, o destinadas a fines distintos de los establecidos para la aplicación de la exención de derechos o de tipo reducido y se hayan abonado los derechos de importación correspondientes. 

El tercero de estos regímenes sería el de \textbf{perfeccionamiento}. Se trata de exenciones aplicadas sobre mercancías extracomunitarias utilizadas en la transformación de mercancías en la Unión o de exenciones aplicadas sobre productos transformados fuera de la Unión con mercancías comunitarias. Como se puede intuir, también existen dos subregímenes: \textbf{perfeccionamiento activo} y \textbf{perfeccionamiento pasivo}. El régimen de Perfeccionamiento activo es un régimen especial mediante el cual, las mercancías no pertenecientes a la Unión pueden ser utilizadas dentro del territorio aduanero de la Unión en una o más operaciones de transformación sin que dichas mercancías, estén sujetas a derechos de importación ni a otro tipo de medidas. Por otro lado, el régimen de Perfeccionamiento pasivo es un régimen especial que permite exportar temporalmente mercancías de la Unión Europea fuera del territorio aduanero de la misma para someterlas a operaciones de transformación y despachar a libre práctica, con exención total o parcial de los derechos de importación, los productos transformados que resulten de esas operaciones al entrar de nuevo en la UE.\footnote{%
    Este último está regulado en el Título VII del Reglamento (UE) nº 952/2013 del Parlamento Europeo y del Consejo, por el que se establece el Código aduanero de la Unión (en adelante CAU), y en el mismo título del Reglamento Delegado (UE) 2015/2446 de la Comisión por el que se completa el CAU y del Reglamento de Ejecución (UE) 2015/2447 de la Comisión por el que se establecen normas de desarrollo.
}\textsuperscript{,}\footnote{%
    Atendiendo a la normativa comunitaria, se entiende por operación de transformación en el perfeccionamiento activo: la manipulación de mercancías, incluidos su montaje o ensamblaje o su incorporación a otras mercancías; la transformación de mercancías; la destrucción de mercancías; la reparación de mercancías, incluidas su restauración y su puesta a punto; el uso de mercancías (regulado por la Unión Europea) que no forman parte del producto transformado, pero que permitan o faciliten la producción de éste, incluso aunque se consuman total o parcialmente en el proceso. Existen también algunas particularidades que debe cumplir el \href{https://sede.agenciatributaria.gob.es/Sede/aduanas/regimenes-especiales/perfeccionamiento-activo-perfeccionamiento-pasivo/requisitos-que-debe-reunir-solicitante-activo.html}{sujeto} y el \href{https://sede.agenciatributaria.gob.es/Sede/aduanas/regimenes-especiales/perfeccionamiento-activo-perfeccionamiento-pasivo/condiciones-necesarias-consecucion-autorizacion-perf-act.html}{proceso}.

Las condiciones para la aplicación del perfeccionamiento pasivo son más amplias, pues aplican sobre el sujeto (\href{https://sede.agenciatributaria.gob.es/Sede/aduanas/regimenes-especiales/perfeccionamiento-activo-perfeccionamiento-pasivo/requisitos-que-debe-reunir-solicitante-pasivo.html}{solicitante}), sobre las objeto (\href{https://sede.agenciatributaria.gob.es/Sede/aduanas/regimenes-especiales/perfeccionamiento-activo-perfeccionamiento-pasivo/mercancias-excluidas-regimen-perfeccionamiento-pasivo.html}{mercancías}) y sobre el procedimiento (\href{https://sede.agenciatributaria.gob.es/Sede/aduanas/regimenes-especiales/perfeccionamiento-activo-perfeccionamiento-pasivo/condiciones-necesarias-consecucion-autorizacion-perf-pas.html}{condiciones necesarias}).
}

El último régimen especial es el de \textbf{tránsito}. Existen a su vez dos tipos de tránsito: el tránsito común y el tránsito de la Unión. El régimen de \textbf{tránsito de la Unión} resulta aplicable a las operaciones de tránsito aduanero entre los Estados miembros de la UE (y Andorra y San Marino) y se utiliza, en general, para amparar la circulación de mercancías no pertenecientes a la Unión sin estar sometidas a los derechos de aduanas y a otros gravámenes que en su caso pudieran corresponder (régimen de tránsito externo) y a la circulación de mercancías de mercancías de la Unión que, entre su punto de partida y su punto de destino en la UE, tienen que atravesar el territorio de un tercer país (régimen de tránsito interno). Por otro lado, el régimen de \textbf{tránsito común} se utiliza para amparar la circulación de mercancías entre los Estados miembros de la UE, los países de la AELC (Islandia, Noruega, Liechtenstein y Suiza), Turquía (desde el 1 de diciembre de 2012), la República de Macedonia del Norte (desde el 1 de julio de 2015), Serbia (desde el 1 de febrero de 2016) y Reino Unido (desde el 1 de enero de 2021).

En este caso el pago de la cuota de IVA correspondiente se realiza en el momento de la liquidación de los derechos arancelarios, para ello la administración aduanera expide el modelo 031, de liquidación de IVA, y la carta de pago correspondiente. El denominado DUA (Documento Único Administrativo), es fundamental ya que es el documento básico para poder deducir las cuotas pagadas en las importaciones de bienes, siendo necesario para poder ejercitar este derecho, que estas cuotas estén debidamente registradas en el Libro Registro de Facturas Recibidas, así podremos deducirlo en nuestra declaración de IVA. Los aranceles que se apliquen no serán deducibles en nuestra actividad.

\subsubsection{Medidas de Salvaguarda (SG)}
Las medidas de salvaguardia antes mencionadas (SG) tienen por objeto proteger temporalmente las industrias de la Unión frente al efecto negativo de los incrementos imprevistos y significativos de las importaciones que causen o puedan causarles un grave perjuicio a dichas industrias y se aplican \textit{erga omnes}. La UE dispone, además, de la posibilidad de adoptar medidas para hacer frente a las prácticas comerciales desleales resultantes de las importaciones objeto de dumping o subvencionadas que causen o puedan causar un perjuicio significativo a la industria de la UE. Se trata de las medidas antidumping y antisubvención que la UE puede aplicar de forma selectiva a mercancías importadas de un determinado país o, incluso, por determinadas empresa de un tercer Estado.

También denominadas derechos compensatorios, fueron incorporadas a la legislación comunitaria por el Reglamento 519/94. No obstante, esta ley comunitaria ha sido modificada varias veces, siendo la más reciente la real izada por el Reglamento 478/2015.

Son medidas defensivas adoptadas por la Comisión cuando la evolución de las importaciones de un producto. las cuales suben de forma repentina e imprevisible, provoque o amenace con provocar un grave perjuicio a los productores comunitarios. Generalmente suponen la aplicación de restricciones cuantitativas a la importación.

Estas medidas suponen una protección extraordinaria y temporal, con el objetivo de conferir a la industria afectada un respiro transitorio para reducir la presión de las importaciones y realizar la reconversión necesaria. Además, no tienen un carácter selectivo. Mientras que las medidas antidumping y antisubvención se toman contra países detenninados, las salvaguardias se aplican a las importaciones de todos los países, independientemente de su procedencia. Además, a diferencia de las medidas antidumping y antisubvención, las denuncias solamente pueden ser interpuestas por las autoridades de los Estados miembros.

Si la Comisión considera que existen indicios suficientes, la apertura de la investigación se publicará en el DOUE y se tendrán en cuenta factores como: el volumen y precio de las importaciones; la repercusión para los productos comunitarios, valorando factores como producción, capacidad, existencias, ventas, cuotas, precio, beneficios, empleo, etc.; cualquier otro factor que puede ser detenninante del perjuicio.

La Comisión deberá adoptar una decisión en un plazo de 9 meses desde el momento en que se inició la investigación. Para adoptar las medidas, deberá presentar al Consejo su propuesta de decisión y éste deberá aprobarlas por mayoría cualificada. Pueden adoptarse medidas de salvaguardia provisionales durante un máximo 200 días si se estima que puede haber un daño dificilmente reparable, y la UE debe compensar a los exportadores perjudicados para que no apliquen medidas de retorsión.

El procedimiento puede concluir con la adopción de medidas de salvaguardia definitivas, la cuales pueden tomar la fonna de un aumento de los derechos arancelarios o a través de restricciones cuantitativas·, que sólo deben aplicarse en la medida necesaria: la duración máxima son cuatro años, prorrogables otros cuatro años adicionales como máximo. Cuando el nivel del contingente es igual o superior a la media de las importaciones de los últimos tres años. 

No podrán aplicarse a un país en desarrollo si la cuota en las importaciones comunitarias es inferior al 3%.

El ejemplo contemporáneo canónico es el sector siderúrgico: la vigilancia previa instaurada en 2016 ante la sobrecapacidad mundial del acero suministró la base informativa sobre la que, tras la desviación de flujos comerciales provocada por los aranceles estadounidenses de la Sección 232, se adoptó la salvaguardia del acero mediante el Reglamento de Ejecución (UE) 2019/159, que estableció contingentes arancelarios por familias de productos. La secuencia vigilancia-salvaguardia del acero ilustra mejor que ninguna otra la función sistémica del régimen. Procedimentalmente, el operador solicita el documento ante la autoridad nacional —en España, la Secretaría de Estado de Comercio—, lo obtiene de forma automática y lo presenta en aduana junto con la declaración; el documento es válido en toda la Unión con independencia del Estado miembro de expedición.

\subsection{Instrumentos aplicables a la exportación}
Conocemos ahora todos los principales intrumentos, regímenes y excepciones aplicables sobre cualquier mercancía importada. Sin duda se trata de una de las materias más relevantes de la PCC, aunque no la única. 

Pasamos a analizar los instrumentos aplicables a las mercancías exportadas en virtud de la PCC. Como es natural, la normativa relativa al AAC no prevé la imposición de derecho de aduana a las exportaciones de mercancías de la Unión a países terceros. De esta forma, el régimen general aplicable a dichas exportaciones se contiene en el \href{https://www.boe.es/buscar/doc.php?id=DOUE-L-2015-80561}{Reglamento (UE) n.º 2015/479}.


\subsubsection{Medidas de Salvaguarda: Escasez y/o compromisos}
Sin embargo, el artículo 36 del TFUE prevé la posibilidad de introducir prohibiciones o restricciones a la importación, exportación o tránsito justificadas por razones: i) de orden público, moralidad y seguridad pública; ii) de protección de la salud y vida de personas y animales; iii) de protección del patrimonio artístico, histórico o arqueológico nacional, y iv) de protección de la propiedad industrial y comercial.

No obstante, este Reglamento n.º 2015/479 permite que la Comisión adopte medidas de salvaguarda en caso de escasez de productos esenciales, o para ponerle remedio, y cuando los intereses de la Unión requieran situación crítica de carestía de productos esenciales, o bien para cumplir con los compromisos internacionales. 

Estas medidas de salvaguardia consistirán en la necesidad de conseguir una autorización de exportación. 




\clearpage
\section{Política Comercial Común: Instrumentos de de defensa de la competencia}
\puntosclave{

}



\subsection{Instrumentos de defensa comercial contra prácticas desleales}
Son instrumentos que se emplean \textbf{para asegurar un comercio leal y defender los intereses de las empresas de la UE, según los acuerdos de la OMC}. Es decir, permiten a la UE defender a sus productores contra importaciones o tratamientos desleales, y contra cambios importantes e inesperados en los flujos comerciales, y que además están aceptados en el seno de la OMC. 

Estas medidas se amparan en los códigos sobre barreras no arancelarias de la OMC, y, consecuentemente, pueden dar lugar a recurso al Órgano de Solución de Diferencias (OSO) de la OMC por parte del país afectado, en caso de desacuerdo sobre las medidas adoptadas.

\subsubsection{Medidas antidumping}
Las medidas antidumping, las más utilizadas y controvertidas en la práctica de la Unión, se dirigen contra los productos importados de países terceros, objeto de dumping, que causen o amenacen causar un perjuicio importante a la industria de la UE. 

El dumping es una distorsión muy peculiar que requiere un requisito concreto: que el perjuicio generado se deba a que el productor extranjero está fijando precios inferiores a los que aplica en su país. Cuando concurre esta circunstancia, que es una práctica prohibida por el artículo VI del GATT, se puede aplicar una contramedida que aumenta los derechos de aduana con el fin de nivelar los precios, práctica que se ha llamado \textbf{medidas anti-dumping}. Tras el desarrollo de un procedimiento administrativo, el Reglamento 2016/1036, la imposición de medidas antidumping exige la concurrencia de las siguientes condiciones:
\begin{lnum}
    \item Las importaciones deben ser objeto de dumping, es decir, que el precio de un producto importado a la UE sea inferior a su valor normal;
    \item Debe haber un perjuicio para la industria de la UE que produce el producto similar. Es decir, que se produzca un aumento de las importaciones a unos precios significativamente inferiores y que afectan a las ventas, rentabilidad, cuota de mercado o productividad de los productores comunitarios;
    \item Debe existir un nexo causal entre las importaciones objeto de dumping y el perjuicio;
    \item La medida antidumping no debe ir en contra de los intereses de la UE, es decir, que el daño causado por las medidas al conjunto de su economía no debe ser mayor que el alivio aportado a la industria que sufre por las importaciones.
\end{lnum}

Cuando concurren estas condiciones, los productores de la UE del producto en cuestión, o en su nombre, presentan una denuncia antidumping ante la Comisión, ya sea directamente o a través de las autoridades de un país de la UE. Los sindicatos también pueden presentar denuncias conjuntamente con la industria de la UE y convertirse en partes interesadas en los procedimientos. En circunstancias especiales, la Comisión también puede abrir una investigación por iniciativa propia.

Toda denuncia debe incluir elementos de prueba de dumping, de perjuicio y del nexo causal entre las importaciones presuntamente objeto de dumping y el supuesto perjuicio. La Comisión debe examinar la exactitud e idoneidad de los elementos de prueba presentados en la denuncia para determinar si existe una base suficiente que justifique la apertura de una investigación. La Comisión deberá iniciar la investigación en el plazo de 45 días a partir de la fecha de presentación de la denuncia.

Una vez que la Comisión decide iniciar una investigación antidumping, debe publicar un anuncio a tal efecto en el Diario Oficial de la Unión Europea (DOUE). Asimismo, se pondrá en contacto con todos los fabricantes conocidos y con todas las demás partes interesadas, pidiéndoles que rellenen los cuestionarios pertinentes en un plazo estricto. 

Cuando existan numerosas partes potencialmente interesadas, la Comisión podrá decidir realizar la investigación sobre la base de una muestra de operadores (productores exportadores, fabricantes de la UE, importadores, usuarios). La Comisión podrá establecer medidas antidumping provisionales transcurridos 60 días desde la fecha de apertura de la investigación, pero antes de los siete meses desde la fecha de inicio del procedimiento. Dichas medidas provisionales tendrán una duración de seis meses, prorrogables hasta tres meses, o se podrán establecer los nueve meses directamente en situaciones excepcionales.

Cuando, en base a la investigación, la Comisión considere que se ha producido un dumping, se podrán imponer medidas antidumping a las importaciones del producto en cuestión en la UE, que generalmente adoptan la forma de un derecho de aduana ad valorem; derechos específicos; o un compromiso de precios en los que el exportador respete los precios mínimos de importación. La investigación hasta la imposición de medidas definitivas será de doce meses, excepcionalmente quince.

En caso de establecerse unas medidas definitivas, éstas podrán permanecer en vigor durante cinco años. Al cabo de ese tiempo, las medidas expirarán a menos que se concluya de que, si la medida expirara, sería probable que continuara o se repitiera el dumping y el perjuicio significativo. Los derechos correrán a cargo del importador de la UE y serán percibidos por las autoridades aduaneras nacionales de los países de la UE afectados. Los importadores podrán solicitar la devolución total o parcial de los derechos pagados si pueden demostrar que el margen de dumping, sobre cuya base se pagaron los derechos, ha sido eliminado o reducido.

En determinadas condiciones, las medidas en vigor podrán reconsiderarse (reconsideración provisional). El alcance de esta reconsideración suele limitarse a uno o varios elementos de las medidas iniciales. Por ejemplo, el nivel de dumping y/o el perjuicio, el ámbito de aplicación del producto o el tipo de medidas.


\subsubsection{Medidas antisubvención}
Las medidas antisubvención o derechos compensatorios son impuestos por la Unión a productos que se han beneficiado en sus países de origen de una subvención, aunque únicamente generarán medidas compensatorias las subvenciones específicas a una empresa o grupo de empresas o industrias, excluyéndose las de investigación o las concedidas a regiones desfavorecidas. Se rigen por el Reglamento 2016/1037. 

En esencia, se busca perseguir una posición económica ventajosa, que sitúa el producto de importación en una posición más favorable sin deberse a las condiciones de mercado.\footnote{%
    Se confiere un beneficio cuando las condiciones en las que se otorga cualquiera de estas contribuciones son más favorables de las que están disponibles en el mercado. Por ejemplo, si un gobierno suministra electricidad por debajo del precio de mercado o compra un producto por encima de su valor de mercado.
} Por tanto, pueden ser el resultado de ayudas económicas directas o indirectas, pero siempre por parte de las autoridades de su país. Las subvenciones pueden revestir diferentes formas: donaciones, préstamos, incentivos fiscales, o bienes o servicios prestados por los poderes públicos. Los aspectos relativos al procedimiento son prácticamente idénticos a las recogidos en el Reglamento 2016/1036 antidumping, salvo que la Comisión debe ofrecer consultas al Gobierno del país exportador antes de iniciar la investigación, que las medidas provisionales duran cuatro meses no prorrogables y que las medidas provisionales se aplican entre los sesenta días y los nueve meses desde el inicio de la investigación.

El derecho compensatorio también se traducirá en un aumento del derecho de aduana aplicable a la mercancía importada, aunque las autoridades pueden también comprometerse a suprimir la subvención o a tomar medidas en relación con los efectos de la misma. 

Tanto las medidas antidumping como las medidas antisubvención tienen que ser compatibles como las normas sobre ambos instrumentos de defensa comercial establecidas por el derecho de la OMC.

\subsection{Intrumentos de Nueva Generación}
Al margen de estos dos instrumentos que se han venido utilizando tradicionalmente para la defensa ante la competencia desleal dentro del marco de la OMC, la Union ha ido desarrollando a lo largo de los últimos veinte años una serie de instrumentos de \textbf{nueva generación} con la intención de adaptarse a la nueva realidad.\footnote{%
    Son el Mecanismo de Control de las Inversiones Extranjeras (Investment Screening), la modificación del Reglamento de Enforcement, el Instrumento de Contratación Pública Internacional (IPI); el Reglamento de Subvenciones Extranjeras (FSR, por sus siglas en inglés); otros se encuentran en fase de aprobación, como son el Instrumento Anti-coerción (ACI, por sus siglas en inglés), el Mecanismo de Ajuste en Frontera por Emisiones de Carbono (CBAM, por sus siglas en inglés) o el Reglamento de Deforestación.
}

Presentamos aquí una pequeña muestra de estos nuevos instrumentos, muchos de ellos presentados en 2021 junto con la estrategia comercial de la Unión (Trade Policy Review) y que persiguen la idea de democratizar las prácticas de solución de conflictos en el ámbito comercial, reducir las asimetrías competitivas (distorsiones, pero no) que generan las preferencias de la Unión y derivan en las exigencias impuestas por su ordenamiento jurídico a los productores locales, y aumentar la autonomía estratégica de la Unión en un mundo cada vez más fragmentado.

\subsubsection{Reglamento de Óbstaculos al Comercio (ROC)}
La UE dispone de otro mecanismo de de defensa comercial entendida en sentido amplio. Se trata del instrumento denominado \textbf{Reglamento de Obstáculos al Comercio (ROC)}, incorporado a la legislación comunitaria con el Reglamento 3286/1994 y modificado en varias ocasiones, la más reciente con el Reglamento 2015/1843.

Tiene como objetivo permitir a empresas, sectores económicos y Estados miembros denunciar obstáculos comerciales en mercados de terceros países. Esto se consigue facilitando el suministro de información la Unión informaciones con el fin de que pueda reaccionar ante los obstáculos al comercio adoptados o mantenidos por terceros países, y contravegan el marco de la OMC. El procedimiento debe iniciarse a instancia de parte presentando una denuncia por escrito ante la Comisión, haciendo constar: 
\begin{lnum}
    \item La especificación de los bienes o servicios afectados por las barreras;
    \item Las pruebas que permitan pensar que existe dicha barrera;
    \item Pruebas de que las barreras comerciales están ocasionando efectos adversos. 
\end{lnum}

Tras esto, la Comisión dispone de $45$ días para decidir si inicia una investigación, y en caso afirmativo será publicado en el DOUE. Durante el proceso, la Comisión recopila información y presenta propuestas a los terceros países involucrados. Esta parte del procedimiento puede durar hasta cinco meses, ampliables hasta siete para casos más complejos. Al finalizar se presenta un informe a los Estados afectados.

En caso de que la investigación concluya la necesidad por parte de la UE de aprobar alguna medida, la Comisión realizará consultas bilaterales con el país infractor. En este caso, se puede conseguir que el país que ha puesto el obstáculo elimine unilateralmente los efectos comerciales adversos o el perjuicio, o que se llegue a un acuerdo bilateral amistoso entre la UE y el tercer país.\footnote{%
    En caso contrario existen varias soluciones. Por un lado, la Unión podrá recurrir al procedimiento de solución de diferencias de la OMC si se vulnera un acuerdo de la OMC. Por otro lado, podrá establecer medidas para corregir el obstáculo (medidas comerciales de retorsión) que serán propuestas por la Comisión y adoptadas por el Consejo en un plazo de $30$ días desde su presentación. Esas medidas de retorsión pueden consistir en: (i) la suspensión o retirada de cualquier concesión derivada de negociaciones de política comercial; (ii) el aumento de derechos de aduana o cualquier otro gravamen a la importación; o, (iii) el establecimiento de restricciones cuantitativas.
}

Como se puede ver, este Reglamento ROC es un paliativo al carácter interestatal del sistema de solución de diferencias de la OMC, al que no tiene acceso los particulares, ya que permite a la Unión activar este mecanismo utilizando la valiosa información que le facilitan las empresas de la UE que padecen las consecuencias de los obstáculos comerciales ilícitos de terceros Estados.

\subsubsection{Instrumento Anticoerción (ACI)}
Se trata de un instrumento jurídico, adoptado en 2023 mediante el \href{https://www.boe.es/buscar/doc.php?id=DOUE-L-2023-81762}{Reglamento 2023/2675}, cuyo objetivo es proteger a los ciudadanos de la Unión de coerciones extranjeras (nacionales están prohibidas de base) en el ámbito comercial. 

La coerción en el contexto comercial se da cuando un país no perteneciente a la Unión intenta presionar a ésta o a uno de sus Estados miembros para que modifique su política en un ámbito concreto, aplicando (o amenazando con aplicar) medidas que afecten al comercio o a la inversión. Se trata de una forma de chantaje que puede afectar a cualquier ámbito político y adoptar múltiples formatos, tanto legislativos como de otro tipo.

¿Cómo funciona el nuevo instrumento anticoercitivo? El instrumento anticoercitivo es la primera herramienta que ha tenido la Unión para tratar con países que restringen el comercio en un intento de forzar un cambio en las políticas europeas. La correspondiente legislación entró en vigor a finales de 2023. El objetivo de este instrumento, al que aún no se ha recurrido, es actuar como medida disuasoria y enviar a los socios comerciales internacionales el mensaje de que la Unión no está dispuesta a tolerar la coerción económica. Le permite señalar formalmente casos de coerción económica por parte de terceros países e intentar el diálogo y la negociación. Después, si la diplomacia convencional fracasa, la Unión puede recurrir, como último recurso, a una amplia gama de contramedidas. Entre las posibles medidas figuran imponer aranceles a productos, restringir servicios, retirar la protección de la propiedad intelectual y la inversión extranjera y eliminar el derecho de presentar ofertas a licitaciones públicas. Estas medidas pueden aplicarse por la vía rápida, pero deben ser proporcionadas y específicas, y solo pueden mantenerse mientras persista la coerción.

Entre los casos más emblemáticos que motivaron su adopción está el de China. China impuso restricciones comerciales a Lituania después de que esta anunciara la mejora de sus relaciones comerciales con Taiwán en junio de 2021. Unos meses después del anuncio, varias empresas lituanas declararon tener dificultades para renovar o formalizar contratos con empresas chinas. Además, tuvieron problemas porque sus envíos no se despachaban y tampoco les era posible presentar la documentación aduanera. El Parlamento ha denunciado la coerción económica de China contra Lituania en varias resoluciones. Más recientemente, \href{https://www.youtube.com/watch?v=lIg4BZe9U-8}{el presidente de Estados Unidos}, Donald Trump, amenazó con imponer aranceles a las mercancías procedentes de una serie de países europeos por su oposición a una posible anexión estadounidense de Groenlandia, territorio autónomo de un país de la UE (Dinamarca).

\subsubsection{Aproximación holística: ¿nuevas bases para la PCC?}
Los hasta ahora presentados se apoyan directamente el art. $207$ TFUE relativo a la Política Comercial Común. No obstante, la Comisión ha adoptado una visión más holística de los instrumentos necesarios para conseguir sus objetivos. De esta forma, algunos de los instrumentos de nueva generación, que ahora veremos, encuentran su fundamento jurídico en otras partes del TFUE, a pesar de la incidencia que tienen en el comercio.

\paragraph{Mercado Interior (\textit{ex} art. $114$ TFUE): Reglamento de Subvenciones Extranjeras (FSR)}
El Reglamento 2022/2560 de Subvenciones Extranjeras (FSR) entró en vigor el 12 de enero de 2023 con el objetivo de hacer frente a las distorsiones originadas por las subvenciones extranjeras en el mercado interior de la UE.

El FSR no debe confundirse con el Reglamento antisubvenciones tradicional (Reglamento (UE) 2016/1037), que se enmarca en la defensa comercial clásica y presenta tres limitaciones que el FSR viene a superar:
\begin{lnum}
    \item \textit{Ámbito material}: el instrumento clásico solo permite imponer derechos compensatorios sobre \textbf{importaciones de mercancías} subvencionadas que causan un perjuicio a la industria de la Unión. El FSR, en cambio, es un instrumento horizontal que cubre \textbf{cualquier actividad económica} realizada en el mercado interior por una empresa que haya recibido financiación extranjera, incluidas las concentraciones (fusiones y adquisiciones), la contratación pública y, con carácter residual, cualquier otra situación de mercado (inversiones, ayudas a la explotación, servicios).
    \item \textit{Test aplicable}: mientras que el mecanismo clásico exige acreditar un perjuicio importante a la industria de la Unión derivado de importaciones subvencionadas, el FSR aplica un test de \textbf{distorsión de la competencia}, más próximo en su construcción al concepto de ayuda de Estado del artículo 107 TFUE (ventaja, imputabilidad a un poder público, selectividad y afectación a la competencia) que a la disciplina antisubvenciones de la OMC.
    \item \textit{Mecanismo de control}. El FSR introduce, además de la investigación de oficio, un régimen de \textbf{notificación previa obligatoria} para concentraciones (a partir de 500 millones de euros de volumen de negocios y 50 millones de contribución financiera extranjera) y para licitaciones públicas (a partir de 250 millones de euros de valor del contrato y 4 millones de contribución financiera extranjera por tercer país), algo ajeno al sistema antisubvenciones tradicional, que opera exclusivamente a instancia de la industria afectada.
\end{lnum}

A diferencia de los instrumentos clásicos de defensa comercial, anclados en el artículo $207$ TFUE (Política Comercial Común), el FSR se basa en el artículo $114$ TFUE, relativo a la aproximación de legislaciones para el establecimiento y funcionamiento del mercado interior. Esta elección no es casual: permite a la Comisión actuar sobre empresas y operaciones que tienen lugar \textit{dentro} de la Unión (y no solo sobre mercancías que cruzan su frontera), evitando así el vacío de protección que dejaba el hecho de que las ayudas de Estado del artículo 107 TFUE solo se aplican a subvenciones concedidas por Estados miembros, y no por terceros países.

El primer caso resuelto con condiciones en materia de contratación pública se produjo el 21 de abril de 2026, cuando la Comisión autorizó la adjudicación de la Línea Violeta del metro de Lisboa condicionada a la sustitución de un subcontratista beneficiario de subvenciones chinas por un fabricante polaco no subvencionado. El 11 de diciembre de 2025 se abrió además una investigación exhaustiva sobre la empresa china Nuctech en el ámbito de los sistemas de detección de amenazas.

\paragraph{Política Medioambiental (\textit{ex} art. $192$): Mecanismo de Ajuste en Frontera de Emisiones de Carbono (CABM)}
El Reglamento (UE) 2023/956, de 10 de mayo de 2023, crea el \textbf{CBAM} (\textit{Carbon Border Adjustment Mechanism}), instrumento que fija un precio al carbono incorporado en determinadas importaciones intensivas en emisiones (hierro, acero, aluminio, cemento, fertilizantes, hidrógeno y electricidad), con el fin de evitar la fuga de carbono hacia jurisdicciones con normativa climática menos exigente y de igualar las condiciones de competencia entre productores europeos e importadores.

Por tanto, el CBAM se apoya en el artículo 192, apartado 1, TFUE (política medioambiental), y no en el artículo 207 TFUE, pese a operar formalmente como un ajuste en frontera con efectos directos sobre el comercio internacional. Esta elección de base jurídica es, precisamente, uno de los puntos más discutidos por la doctrina.

El Reglamento contempló un período transitorio (octubre de 2023 a diciembre de 2025) de mera obligación informativa, sin pago alguno. La \textbf{fase definitiva} comenzó el 1 de enero de 2026: desde entonces, solo pueden importarse mercancías CBAM a través de declarantes autorizados y registrados, y los importadores deberán adquirir certificados CBAM equivalentes a las emisiones declaradas, cuyo precio se referencia al del régimen de comercio de derechos de emisión (ETS). El Reglamento (UE) 2025/2083, adoptado el 8 de octubre de 2025 dentro del primer paquete de simplificación (\textit{Omnibus}), introdujo además un umbral de minimis de $50$ toneladas anuales por importador, por debajo del cual no existen obligaciones CBAM, pensado sobre todo para aliviar a las PYMES.

El análisis del \href{https://www.realinstitutoelcano.org/analisis/el-arancel-al-carbono-cbam-proteccionismo-verde-o-liderazgo-global-contra-el-cambio-climatico/}{Real Instituto Elcano} (FEÁS et al.) sostiene que el CBAM no es, en su diseño, un mecanismo proteccionista ni recaudatorio, sino un instrumento para sostener el nivel de ambición climática de la Unión sin penalizar a su industria; sin embargo, reconoce que será difícil que la UE evite las acusaciones de proteccionismo verde y unilateralismo por parte de terceros países, y anticipa que es probable que economías emergentes impugnen la medida ante el sistema de solución de diferencias de la OMC.\footnote{%
    Desde la perspectiva del Derecho de la OMC, la compatibilidad del CBAM con el GATT dependerá de si logra ampararse en las excepciones generales del artículo XX (apartados b y g, relativos a la protección de la salud y a la conservación de recursos naturales agotables), lo que exige superar el exigente test del \textit{chapeau} --no discriminación arbitraria o injustificable-- fijado por la jurisprudencia del Órgano de Apelación (\textit{Estados Unidos -- Camarones}).
} El debate de fondo, aún no resuelto, es si la magnitud real del riesgo de fuga de carbono justifica una medida unilateral de este alcance frente a la vía multilateral.

\paragraph{Política Medioambiental (\textit{ex} art. $192$): Reglamento de Deforestación (EUDR)}
El Reglamento (UE) 2023/1115, de 31 de mayo de 2023 (\textbf{EUDR}, \textit{EU Deforestation Regulation}), entró en vigor el 29 de junio de 2023 y establece obligaciones de diligencia debida para los operadores y comerciantes que introduzcan, comercialicen o exporten desde la Unión determinadas materias primas asociadas a la deforestación (ganado vacuno, cacao, café, aceite de palma, caucho, soja y madera) y sus productos derivados (cuero, chocolate, muebles, entre otros). Solo pueden comercializarse en la UE los productos que estén libres de deforestación posterior al 31 de diciembre de 2020, que cumplan la legislación del país de producción y que vayan acompañados de una declaración de diligencia debida con las coordenadas geográficas de las parcelas de origen.

Como se ve, al igual que el CBAM, el EUDR se sustenta en el artículo $192$ TFUE (política medioambiental), en este caso vinculado a los objetivos de la Estrategia de la UE sobre Biodiversidad para 2030, y no en la Política Comercial Común, pese a condicionar directamente el acceso al mercado interior de mercancías procedentes de terceros países.

A diferencia de lo previsto originalmente, el Reglamento \textbf{no ha entrado aún en aplicación} y acumula ya dos aplazamientos. La fecha de aplicación inicial (30 de diciembre de 2024) se pospuso un año mediante el Reglamento (UE) 2024/3234, hasta el 30 de diciembre de 2025 para grandes y medianos operadores (30 de junio de 2026 para micro y pequeñas empresas). Un segundo aplazamiento, aprobado mediante el Reglamento (UE) 2025/2650, de 19 de diciembre de 2025, retrasa la aplicación hasta el 30 de diciembre de 2026 para grandes y medianos operadores, y hasta el 30 de junio de 2027 para micro y pequeñas empresas, con el fin de dar más tiempo a operadores y autoridades para adaptarse al sistema informático de trazabilidad. El EUDR derogará el Reglamento de la Madera (EUTR), aunque este seguirá aplicándose transitoriamente hasta el 31 de diciembre de 2029 respecto de los productos de madera producidos con anterioridad a la entrada en vigor del EUDR.

La doctrina señala dos frentes críticos recurrentes. Por un lado, la capacidad real de pequeños productores de países en desarrollo para cumplir con la exigencia de geolocalización parcela por parcela, lo que genera un riesgo de exclusión de facto de operadores que no disponen de medios técnicos ni administrativos suficientes. Por otro lado, la tensión entre el EUDR y la política comercial exterior de la Unión, singularmente visible en las negociaciones del acuerdo UE-Mercosur, donde exportadores como Brasil han cuestionado la compatibilidad de exigencias medioambientales unilaterales con los compromisos de apertura comercial recíproca. Los sucesivos aplazamientos --motivados oficialmente por la falta de preparación técnica del sistema de información TRACES-NT-- son también leídos por parte de la doctrina como una cesión política frente a la presión de terceros países y de sectores productivos europeos, lo que debilita la credibilidad del instrumento como pieza del Pacto Verde Europeo.





\clearpage
\section{Política Comercial Común: Regímenes comerciales preferenciales}
\puntosclave{
    La política comercial convencional es aquella que se define mediante la red de acuerdos bilaterales y plurilaterales en el marco de la OMC.
    
    A nivel multilateral, el fracaso de la Ronda Doha y de la política del single undertaking ha llevado a que, desde la Conferencia Ministerial de Bali (2013), se opte por una política de early harvest. La XII Conferencia Ministerial de Ginebra (2022) ha sido la más exitosa desde la creación de la OMC en 1995, principalmente por el acuerdo sobre pesca ilegal. 
    
    La UE forma parte de los tres acuerdos plurilaterales existentes el marco de la OMC: TPA, ITA y GPA. 
    
    Los acuerdos de «Nueva Generación» que firma la UE negocia a nivel birregional y bilateral se caracterizan por su amplio contenido (ya no solo el relativo a las 111ercancías, como en los acuerdos de «Primera Generación»), por la consideración del desarrollo sostenible (TSD), y por la introducción de cuestiones relativas a la gobemanza y resolución de disputas en materia de inversión.
}


Hasta ahora hemos presentado todos los instrumentos autónomos de la Política Comercial. Hoy en día, siendo la OMC una Organización multilateral que agrupa a la gran mayoría de Estados, lo hasta ahora presentado sería \textit{los instrumentos autónomos dentro del régimen multilateral de la Política Comercial de la Unión}. 

Ahora bien, existe una gran variedad de compromisos, adoptados mediante Acuerdo, a través de los cuales se formalizan los \textbf{regímenes comerciales preferenciales}, de gran varieda, importancia y con prácticamente todos los países del mundo. Este régimen especial es lo que hoy en día conviene entender como la Política Comercial convencional, experimentando de esta forma una evolución doctrinal más adecuada al contexto actual. 

Para presentar esta vertiente tan importante de la Política Comercial Común, cuyo fundamento jurídico conviene repetir que se encuentra en el art. $209$ TFUE (capacidad para adoptar Acuerdos de Asociación), presentaremos en primer lugar los instrumentos tradicionales y en segundo lugar lo que se conoce como Acuerdos Comerciales de Nueva Generación.

\subsection{Marco normativo (\textit{ex} arts. $217$ y $218$ TFUE): Fundamentos y procedimiento}
La celebración de regímenes comerciales preferenciales se apoya en un entramado de tres disposiciones del TFUE que conviene distinguir con precisión, puesto que cada una cumple una función distinta y su confusión es un error frecuente. El primero, art. $207$ TFUE que constituye la base material de la Política Comercial Común, ya lo tratamos en el primer apartado y vimos como define su ámbito de aplicación.

El segundo, e importante para lo que nos ocupa, es el artículo $217$ TFUE que habilita a la Unión para concluir con uno o varios terceros países u organizaciones internacionales acuerdos por los que se cree una asociación que suponga derechos y obligaciones recíprocos, acciones comunes y procedimientos particulares, constituye la base jurídica de los llamados Acuerdos de Asociación y se distinguen así de los simples acuerdos comerciales celebrados exclusivamente sobre la base del artículo 207 TFUE, sin marco de asociación política. 

El artículo $218$ TFUE regula, finalmente, el procedimiento de negociación y celebración de todos estos acuerdos, con independencia de que su base material sea el artículo $207$ o el $217$ TFUE. El procedimiento se articula en las siguientes fases: (i) la Comisión presenta una recomendación al Consejo; (ii) el Consejo autoriza la apertura de las negociaciones y adopta las directrices de negociación, pudiendo designar un comité especial al que la Comisión debe consultar a lo largo del proceso (Comité de Política Comercial); (iii) concluida la negociación, el Consejo, a propuesta del negociador, adopta la decisión que autoriza la firma y, en su caso, la aplicación provisional del acuerdo; y, (iv) el Consejo adopta la decisión de celebración del acuerdo, que en la generalidad de los acuerdos comerciales y, en todo caso, en los Acuerdos de Asociación, requiere la previa aprobación (dictamen conforme) del Parlamento Europeo, conforme al artículo 218.6.a) TFUE.\footnote{%
    Este es uno de los grandes cambios introducidos por el Tratado de Lisboa: antes de 2009, el Parlamento Europeo solo era consultado en materia comercial; desde entonces, su aprobación es preceptiva para la práctica totalidad de los acuerdos comerciales, lo que ha reforzado sustancialmente su papel en la PCC.
} 

En cuanto a la votación en el seno del Consejo, la regla general es la mayoría cualificada, pero el TFUE (art. $218.8$) exige unanimidad para los Acuerdos de Asociación, para los acuerdos en todos los ámbitos en que la adopción de un acto interno de la Unión requeriría unanimidad.

\subsubsection{El problema de los acuerdos mixtos}
Buena parte de los Acuerdos de Asociación y algunos ALC de Nueva Generación no se agotan en materias de competencia exclusiva de la Unión, sino que incorporan también disposiciones sobre materias de competencia compartida entre la Unión y los Estados miembros. Típicamente, la inversión extranjera distinta de la directa (inversión de cartera) y el mecanismo de solución de controversias inversor-Estado, tal y como quedó fijado por el TJUE en el ya citado Dictamen 2/15. 

Cuando esto ocurre, el acuerdo se califica como \textbf{acuerdo mixto}, lo que tiene una consecuencia procedimental de primer orden: además de la aprobación por las instituciones de la Unión conforme al procedimiento del artículo $218$ TFUE, resulta necesaria la ratificación individual por parte de los veintisiete Estados miembros, conforme a sus respectivos procedimientos constitucionales internos, lo que en algunos Estados miembros implica, además, la intervención de parlamentos regionales o subestatales.

Esta exigencia adicional de ratificación nacional (y, en su caso, regional) es la causa de los prolongados retrasos que sufren muchos de estos acuerdos entre su firma y su entrada en vigor plena. El episodio más conocido es el bloqueo temporal que el Parlamento regional de Valonia impuso a la firma del CETA en octubre de 2016, en ejercicio de la competencia que la Constitución belga reconoce a sus entidades federadas sobre la celebración de tratados internacionales, y que solo se resolvió mediante garantías adicionales incorporadas a una declaración interpretativa conjunta. 

El caso de Mercosur, ya mencionado, ilustra una variante distinta del mismo problema estructural: en enero de 2026 el Parlamento Europeo decidió remitir el Acuerdo de Asociación al TJUE para que se pronuncie sobre su compatibilidad con los Tratados antes de continuar con el procedimiento de ratificación, lo que ha paralizado de facto su entrada en vigor plena. Si bien el Acuerdo Interino de Comercio, al no requerir ratificación de los Estados miembros por limitarse a las materias de competencia exclusiva de la Unión, sí ha podido aplicarse provisionalmente desde mayo de 2026.\footnote{%
    Esta distinción, Acuerdo de Asociación mixto frente a Acuerdo Interino de Comercio de competencia exclusiva, es precisamente la técnica que la Comisión ha generalizado desde el caso CETA para blindar la aplicación de la parte estrictamente comercial de estos acuerdos frente a los riesgos de bloqueo en la ratificación nacional.
}

\subsubsection{Taxonomía: la escala de intensidad de las preferencias comerciales}
Sentado lo anterior, los distintos regímenes comerciales preferenciales pueden ordenarse, siguiendo el criterio doctrinal español (MANGAS MARTÍN) de la intensidad de la preferencia concedida, en una escala descendente que responde, a su vez, a una lógica política subyacente: cuanto mayor es la proximidad geográfica, la perspectiva de adhesión o la afinidad de valores con el socio, mayor es la profundidad de integración que la Unión está dispuesta a conceder; cuanto menor es esa proximidad, o mayor la asimetría de desarrollo económico, más se desplaza el acuerdo hacia fórmulas de cooperación o hacia la concesión unilateral:
\begin{lnum}
    \item \textbf{Regímenes de integración máxima}. Suponen la extensión, parcial o total, del propio mercado interior de la Unión (o de la unión aduanera) a un tercer país, sin que este llegue a ser miembro.
    \item \textbf{Acuerdos comerciales de «Nueva Generación»}. ALC recíprocos de contenido muy amplio (mercancías, servicios, inversiones, contratación pública, propiedad intelectual, TSD y gobernanza reforzada) con economías desarrolladas o emergentes de peso.
    \item \textbf{Acuerdos de segunda línea}. ALC recíprocos igualmente, pero de contenido y ambición más limitados.
    \item \textbf{Acuerdos de Asociación Económica (EPA) con países ACP}. En los que el componente de cooperación al desarrollo prevalece sobre la reciprocidad comercial plena.
    \item \textbf{Sistema de Preferencias Generalizadas (SPG)}. Único régimen de base autónoma y unilateral, sin fundamento convencional, en el que el tercer país beneficiario no negocia ninguna contrapartida.
\end{lnum}

Cada uno de estos niveles se desarrolla en los epígrafes siguientes, dedicando un apartado propio a cada acuerdo de relevancia dentro de cada nivel.


\subsection{Regímenes de integración máxima}
Los tres regímenes de integración máxima comparten el rasgo de conceder al tercer país un grado de integración económica con la Unión sensiblemente superior al de cualquier acuerdo de libre comercio, pero difieren radicalmente en su naturaleza jurídica, en su alcance material y en su función. Ilustran, en definitiva, tres soluciones distintas a un mismo dilema: \textbf{cómo conceder un grado de integración económica muy elevado a un tercer país sin extenderle la condición de Estado miembro}.

El Espacio Económico Europeo constituye una extensión del mercado interior de la Unión --con sus cuatro libertades-- a Estados que han optado deliberadamente por no ser miembros; la unión aduanera con Turquía constituye, en cambio, la integración en el territorio aduanero único de la Unión de un Estado que sí aspira a la adhesión, pero limitada al comercio de mercancías; y los Acuerdos de Estabilización y Asociación con los Balcanes Occidentales no son, en rigor, un régimen de integración económica autónomo, sino el instrumento jurídico de preadhesión --basado, junto con el resto, en el artículo 217 TFUE-- mediante el que la Unión acompaña a un grupo de países candidatos o candidatos potenciales en su proceso de aproximación al acervo comunitario, con vistas a una futura adhesión plena. Los tres regímenes 

\subsubsection{Espacio Económico Europeo (EEE, 1992)}
El Acuerdo sobre el Espacio Económico Europeo, firmado en Oporto en mayo de 1992 y en vigor desde el 1 de enero de 1994, extiende a Islandia, Liechtenstein y Noruega (Estados de la Asociación Europea de Libre Comercio, EFTA) las cuatro libertades del mercado interior de la Unión: libre circulación de mercancías, personas, servicios y capitales. Suiza, cuarto miembro de la AELC, rechazó su participación en el EEE mediante referéndum en diciembre de 1992, y mantiene en su lugar una red propia de acuerdos bilaterales sectoriales con la Unión, estructuralmente distinta y más fragmentada.

Sin embargo, no alcanza la Política Agrícola Común ni a la Política Pesquera Común, que quedan fuera del Acuerdo (si bien existen protocolos específicos sobre determinados productos agrícolas transformados y sobre acceso recíproco a caladeros). Tampoco constituye una unión aduanera ya que los Estados del EEE conservan su propia política comercial exterior y sus propios aranceles frente a terceros países, lo que obliga a mantener normas de origen y controles aduaneros en las fronteras entre el EEE y la Unión.\footnote{%
    Esta es la diferencia estructural más relevante frente a la unión aduanera con Turquía: el EEE integra el comercio de mercancías, servicios, personas y capitales, pero no la política comercial exterior; la unión aduanera con Turquía hace exactamente lo contrario.
}

\paragraph{Arquitectura institucional: Dos pilares y principio de homogeneidad} 
El funcionamiento del EEE descansa sobre el llamado \textit{principio de homogeneidad}: el acervo de la Unión relevante para el mercado interior se incorpora de manera dinámica al ordenamiento de los Estados del EEE a través de decisiones del Comité Mixto del EEE, sin que estos Estados participen en el proceso legislativo de la Unión que da origen a esas normas. 

Esta asimetría (los Estados del EEE deben aplicar una normativa en cuya elaboración no han votado) ha sido descrita en la doctrina como el precio de la integración sin membresía, y es objeto de crítica recurrente por parte de sectores euroescépticos noruegos e islandeses, que la han calificado gráficamente de \textit{democracia por fax}. Para preservar la autonomía del ordenamiento jurídico de la Unión, el EEE se articula sobre una estructura institucional de dos pilares: la vigilancia del cumplimiento del Acuerdo por parte de los Estados del EEE corresponde no a la Comisión Europea, sino al Órgano de Vigilancia de la AELC (ESA, por sus siglas en inglés), y el control judicial no corresponde al TJUE, sino al Tribunal de la AELC, que sigue no obstante muy de cerca la jurisprudencia del TJUE para preservar la homogeneidad de interpretación entre ambos pilares. 

Como contrapartida financiera a los beneficios de acceso al mercado interior, Islandia, Liechtenstein y Noruega contribuyen a la cohesión económica y social de la Unión a través de los mecanismos financieros del EEE y de Noruega (\textit{EEA and Norway Grants}).

\subsubsection{Unión Aduanera con Turquía: Acuerdo de Ankara (1963, 1995)}
La relación de asociación entre la Unión y Turquía se remonta al Acuerdo de Ankara de 1963, que estableció un vínculo de asociación con vistas a una futura unión aduanera y, en última instancia, a la adhesión; y al Protocolo Adicional de 1970, que fijó el calendario de desarme arancelario. La fase final de esta unión aduanera se estableció mediante la Decisión n.º 1/95 del Consejo de Asociación CE-Turquía, en vigor desde el 31 de diciembre de 1995.

A diferencia del EEE, que cubre las cuatro libertades pero no la unión aduanera, el régimen turco es exactamente el contrario: constituye una auténtica unión aduanera, con AEC y libre circulación sin controles aduaneros en frontera pero limitada a los productos industriales y a los productos agrícolas transformados, mientras que no participa, en consecuencia, en el mercado interior de servicios ni en las libertades de personas y capitales en los términos en que sí lo hacen los Estados del EEE.

\paragraph{Asimetría estructural: \textit{free rider}} 
El rasgo más criticado por la doctrina es la asimetría que genera la obligación turca de alinear su arancel exterior y su política comercial con los de la Unión para los productos cubiertos, sin que Turquía participe en la negociación de los acuerdos comerciales que la Unión suscribe con terceros países. 

Cuando la Unión concede a un tercer país acceso preferencial a su mercado mediante un ALC, ese tercer país obtiene automáticamente acceso al mercado turco, sin que Turquía obtenga a cambio, salvo que negocie separadamente su propio acuerdo con ese mismo tercer país, un acceso recíproco al mercado de este. Este desequilibrio ha sido uno de los principales argumentos turcos a favor de la modernización del régimen, de manera que Turquía sea consultada o, cuando menos, negocie en paralelo con los socios de los nuevos ALC de la Unión.

Con este fin, la Comisión adoptó ya en diciembre de 2016 una Recomendación al Consejo para negociar directrices de modernización de la unión aduanera (también para extender a servicios, contratación pública y agricultura), pero el proceso quedó congelado por el deterioro de las relaciones políticas y el estancamiento del proceso de adhesión turco desde 2018. La relación se ha reactivado parcialmente en los últimos años: en la reunión celebrada en Ankara en febrero de 2026 entre la comisaria europea de Ampliación, Marta Kos, y el ministro de Asuntos Exteriores turco, Hakan Fidan, ambas partes reafirmaron la voluntad de avanzar hacia la modernización de la unión aduanera, sin que a la fecha se haya concluido un nuevo marco negociador formal.\footnote{%
    Comunicado conjunto UE-Turquía, Ankara, 6 de febrero de 2026, recogido por EFE e Infobae.
} 

Téngase en cuenta que Turquía mantiene el estatuto de país candidato a la adhesión desde 2004, si bien las negociaciones de adhesión propiamente dichas están de facto congeladas desde 2018 por las mismas razones.

\subsubsection{Acuerdos de Estabilización y Asociación (AEA): Candidatos a la Adhesión}
El Proceso de Estabilización y Asociación, lanzado a partir del compromiso político asumido en la Cumbre de Tesalónica de 2003, se instrumenta jurídicamente a través de Acuerdos de Estabilización y Asociación (AEA) bilaterales con cada uno de los socios de la región. A fecha de hoy tienen un AEA en vigor Albania, Bosnia y Herzegovina, Kosovo, Macedonia del Norte, Montenegro y Serbia; el de Kosovo, firmado en 2015 y en vigor desde abril de 2016, fue el primer acuerdo internacional suscrito por la Unión tras la entrada en vigor del Tratado de Lisboa, que le confirió personalidad jurídica propia.\footnote{%
    Esto es particularmente interesante desde una perspectiva política: La Unión firmó este último Tratado en su propio nombre, pues hay países de la Unión (como España) que no reconocen su condición de Estado soberano.
}

A diferencia del EEE y de la unión aduanera turca, que son regímenes concebidos como estables en el tiempo, los AEA son por definición transitorios: su función es preparar progresivamente al socio para la adhesión, y no ofrecerle una alternativa permanente a esta. Por tanto, su contenido habitual combina tres elementos: el establecimiento gradual de una zona de libre comercio bilateral con la Unión, el diálogo político estructurado, y, sobre todo, un compromiso de aproximación progresiva de la legislación nacional al acervo comunitario, sujeto a seguimiento y condicionalidad por parte de la Comisión a través de los sucesivos informes anuales de progreso.

Ahora bien, no todos los socios del Proceso de Estabilización y Asociación tienen la misma condición formal ni el mismo grado de avance. Turquía (2004), Macedonia del Norte (2005), Montenegro (2010), Serbia (2012), Albania (2014) y, más recientemente, Bosnia y Herzegovina (2022) ostentan el estatuto de país candidato oficial; Kosovo, en cambio, permanece como candidato potencial únicamente, al no estar su independencia reconocida por cinco Estados miembros de la Unión (España, Chipre, Grecia, Rumanía y Eslovaquia). El grado de avance en la negociación de capítulos del acervo es también muy dispar: Montenegro encabeza el proceso, habiendo cerrado varios capítulos de negociación en 2025 y con el objetivo declarado de concluir las negociaciones de adhesión a finales de 2026, mientras que Serbia no ha abierto ningún capítulo nuevo en los últimos cuatro años, y Bosnia y Herzegovina arrastra, además, obstáculos de naturaleza constitucional derivados del reparto étnico de cargos públicos declarado contrario al Convenio Europeo de Derechos Humanos por el TEDH en el asunto \textit{Sejdić y Finci c. Bosnia y Herzegovina} (2009).

No obstante, la invasión rusa de Ucrania en febrero de 2022 transformó la política de ampliación hacia los Balcanes Occidentales, que pasó de percibirse como un expediente técnico de bajo interés político a considerarse una cuestión de seguridad estratégica, ante el riesgo de que el vacío de influencia europea en la región fuera ocupado por Rusia o por China (que han intensificado sustancialmente su presencia inversora, particularmente en Serbia). En este contexto, la Comisión propuso en 2023 un Plan de Crecimiento para los Balcanes Occidentales que permite a los socios acceder anticipadamente a determinados beneficios del mercado único y a financiación adicional antes de la adhesión formal, condicionado al cumplimiento de reformas, y ha intensificado la celebración de cumbres UE-Balcanes Occidentales, la más reciente, en Tivat (Montenegro), en junio de 2026.\footnote{%
    Consejo Europeo e Infobae, cobertura de la Cumbre UE-Balcanes Occidentales, Tivat, junio de 2026.
}



\subsection{Acuerdos comerciales de \textit{Nueva Generación}}
Los acuerdos comerciales de Nueva Generación comparten un rasgo definitorio: no se limitan a la reducción de aranceles sobre el comercio de mercancías (propio de los acuerdos de \textit{Primera Generación}), sino que incorporan de manera sistemática capítulos sobre comercio de servicios, inversiones, contratación pública, propiedad intelectual, competencia, cooperación regulatoria y comercio y desarrollo sostenible (TSD), además de un sistema de gobernanza y solución de controversias propio. Su origen se remonta a la estrategia \textit{Global Europe} de la Comisión (2006), formulada como respuesta al estancamiento de la Ronda Doha: ante la imposibilidad de avanzar en la liberalización multilateral, la Unión decidió trasladar su ambición negociadora al plano bilateral y birregional con los socios comerciales de mayor peso económico.

\subsubsection{Bilaterales}

\paragraph{ALC (2011): Unión-Corea del Sur}
El ALC con Corea del Sur, firmado en 2010, fue el primero de esta nueva generación y, en el momento de su celebración, el acuerdo comercial más ambicioso jamás negociado por la Unión en términos de eliminación arancelaria. 

Se aplicó provisionalmente desde el 1 de julio de 2011 y no entró en vigor con carácter definitivo hasta el 13 de diciembre de 2015, una vez concluida la ratificación por la totalidad de los Estados miembros, al tratarse de un acuerdo mixto conforme al modelo anterior al Dictamen 2/15.

\paragraph{El Acuerdo Económico y Comercial Global con Canadá (CETA)}
El CETA, firmado en octubre de 2016 tras siete años de negociación, se aplica provisionalmente desde el 21 de septiembre de 2017, exclusivamente en aquellas materias que son competencia exclusiva de la Unión; el capítulo de protección de inversiones y el Sistema de Tribunales de Inversiones (ICS) quedan excluidos de esa aplicación provisional, precisamente por tratarse de materias de competencia compartida, y su entrada en vigor definitiva sigue pendiente, a la espera de que la ratifiquen todos los Estados miembros.

¿Por qué? La firma del CETA estuvo a punto de frustrarse en octubre de 2016 cuando el Parlamento regional de Valonia, en ejercicio de la competencia que la Constitución federal belga reconoce a sus entidades federadas sobre la celebración de tratados internacionales, se negó a autorizar al Gobierno federal belga a firmar el acuerdo. El bloqueo solo se resolvió mediante la incorporación de un instrumento interpretativo conjunto y de declaraciones adicionales sobre la preservación del poder regulatorio de los entes subestatales, lo que permitió desbloquear la firma pocos días después. Este episodio se convirtió en el ejemplo paradigmático del riesgo que entraña la naturaleza mixta de estos acuerdos.

El CETA introdujo, por primera vez en la práctica convencional de la Unión, un sistema de lista negativa en materia de servicios (todo liberalizado salvo lo expresamente excluido), con compromisos de \textit{standstill} y de apertura futura mediante cláusula \textit{ratchet}; y el ya mencionado Sistema de Tribunales de Inversiones (ICS), primer intento de sustituir el arbitraje ad hoc del ISDS clásico por un órgano jurisdiccional permanente, con vocación de convertirse en el modelo de referencia. De hecho replicado, con matices, en los acuerdos posteriores con Singapur, Vietnam y México.

\paragraph{Nuevo Modelo: ALCs y el APIs (Singapur y Viet Nam)}
El acuerdo con Singapur es el primero en aplicar el modelo escindido derivado del Dictamen 2/15, del que además constituye el origen directo: fue precisamente la consulta planteada por la Comisión sobre la naturaleza de la competencia para celebrar este acuerdo la que dio lugar al Dictamen. Como resultado, se dividió en dos instrumentos: el Acuerdo de Libre Comercio UE-Singapur (EUSFTA), de competencia exclusiva de la Unión, en vigor desde el 21 de noviembre de 2019 sin necesidad de ratificación por los Estados miembros; y el Acuerdo de Protección de las Inversiones UE-Singapur (EUSIPA), de competencia compartida, que sigue pendiente de ratificación por parte de varios Estados miembros.

Siguiendo el mismo modelo escindido inaugurado con Singapur, el Acuerdo de Libre Comercio UE-Vietnam (EVFTA) entró en vigor el 1 de agosto de 2020, mientras que el Acuerdo de Protección de las Inversiones UE-Vietnam (EVIPA) permanece, igualmente, a la espera de la ratificación de los Estados miembros. El EVFTA reviste especial interés porque constituye el primer ALC de Nueva Generación celebrado por la Unión con un país en desarrollo de renta media, lo que se traduce en un capítulo TSD con mayor peso relativo de la asistencia técnica y los periodos transitorios de desarme arancelario, frente al enfoque de reciprocidad plena e inmediata propio de los acuerdos con economías desarrolladas.

\paragraph{Acuerdo de Asociación Económica con Japón}
El Acuerdo de Asociación Económica (EPA) con Japón, en vigor desde el 1 de febrero de 2019, fue en el momento de su celebración el mayor acuerdo de libre comercio bilateral jamás negociado por la Unión, cubriendo entre ambas partes cerca de un tercio del PIB mundial. Se negoció en paralelo, y se firmó el mismo día, un Acuerdo de Asociación Estratégica (SPA) centrado en el diálogo político y la cooperación (separado del EPA precisamente para no supeditar la entrada en vigor de este último a la ratificación mixta que exige el componente político).

No obstante, a diferencia de Singapur y Vietnam, en el caso de Japón las partes no lograron alcanzar consenso sobre el mecanismo de solución de controversias en materia de inversiones, fundamentalmente por la resistencia japonesa a aceptar el modelo ICS defendido por la Unión frente al ISDS clásico. El capítulo de protección de inversiones quedó, en consecuencia, excluido del EPA, cuyo ámbito se limita a la liberalización comercial en sentido estricto. Las negociaciones sobre inversión han continuado por separado, sin que a la fecha se haya alcanzado un acuerdo.

\paragraph{Acuerdo de Comercio y Cooperación con el Reino Unido (TCA, 2021)}
El Acuerdo de Comercio y Cooperación (TCA) con el Reino Unido, aplicado provisionalmente desde el 1 de enero de 2021 y en vigor con carácter definitivo desde el 1 de mayo de 2021, ocupa una posición singular dentro de esta categoría: no responde a una lógica de integración creciente entre dos economías que se aproximan, como el resto de acuerdos de Nueva Generación, sino a la ordenación de la salida de un antiguo Estado miembro del mercado interior.

El Acuerdo garantiza la ausencia de aranceles y de contingentes para el comercio de mercancías originarias de las partes, pero sujeta ese acceso preferencial al cumplimiento de estrictas normas de origen y, sobre todo, a un conjunto de cláusulas de \textit{level playing field} (no regresión en materia social, laboral, medioambiental y de ayudas de Estado) reforzadas por un mecanismo de reequilibrio (\textit{rebalancing mechanism}) que permite a cualquiera de las partes adoptar medidas arancelarias correctoras, unilateralmente, si constata divergencias regulatorias significativas de la otra parte que afecten sustancialmente al comercio o la inversión bilaterales. 

Este mecanismo, sin precedente en los acuerdos de Nueva Generación anteriores, es una de las singularidades más preguntadas del TCA. La gobernanza del acuerdo corresponde a un Consejo de Asociación y varios comités especializados, y su solución de controversias se articula mediante paneles arbitrales ad hoc.


\paragraph{Acuerdo Marco Avanzado con Chile}
El Acuerdo de Asociación UE-Chile, en vigor desde 2003, incluía ya en su origen un componente comercial relativamente ambicioso para su época --liberalización completa de mercancías no agrícolas, cierta apertura de servicios y contratación pública--, gracias al cual el comercio bilateral creció más de un 160\,\% desde su entrada en vigor. Su modernización, negociada entre 2017 y 2022 y concluida políticamente el 9 de diciembre de 2022, ha elevado sin embargo su contenido hasta un nivel que lo sitúa junto a Singapur y Vietnam como referencia del modelo escindido post-Dictamen 2/15, y no junto a los acuerdos andinos.

Siguiendo el modelo ya consolidado, la modernización se articula en dos instrumentos firmados el 13 de diciembre de 2023: el Acuerdo Marco Avanzado (AMA), mixto --pilar político y de cooperación más pilar de comercio e inversión, incluida la protección de inversiones--, sujeto a ratificación por los veintisiete Estados miembros; y un Acuerdo Comercial Interino (ACI), de competencia exclusiva de la Unión, que recibió el consentimiento del Parlamento Europeo el 29 de febrero de 2024 y entró en vigor, tras completar Chile sus propios trámites internos, el 1 de febrero de 2025. El propio pilar político y de cooperación del AMA --con exclusión de las disposiciones comerciales-- se aplica igualmente desde el 1 de junio de 2025, en tanto no concluya la ratificación íntegra del Acuerdo por parte de los Estados miembros.\footnote{%
    EUR-Lex, "Acuerdo Marco Avanzado Unión Europea-Chile"; SUBREI, "Congreso aprueba y despacha Acuerdo Marco Avanzado entre Chile y Unión Europea", noviembre de 2024.
}

El acuerdo modernizado eleva la liberalización de mercancías hasta el 99,9\,\% de las exportaciones de la Unión a Chile libres de arancel --manteniéndose exclusiones puntuales para el azúcar en ambas partes, y para la Unión en plátanos y arroz--, protege 216 indicaciones geográficas europeas (43 de ellas españolas) y mejora sustancialmente el acceso recíproco a la contratación pública, incluidos, por primera vez, los contratos relacionados con materias primas críticas y energía.

El capítulo 17 del AMA crea un Sistema de Tribunales de Inversiones (ICS) de sede permanente --con un tribunal de primera instancia y un tribunal de apelación, integrados por miembros independientes--, que sustituye y deroga los dieciséis Acuerdos de Protección Recíproca de Inversiones bilaterales que Chile mantenía con distintos Estados miembros, unificando así un régimen hasta entonces fragmentado y heterogéneo.\footnote{%
    Biblioteca del Congreso Nacional de Chile, Informe sobre el AMA UE-Chile, 2025; ISDS América Latina, "Chile", 2024.
}

El AMA con Chile es, además, el primer acuerdo económico-comercial de la Unión que incorpora un capítulo específico de género y comercio, orientado a promover la igualdad de oportunidades y de trato entre hombres y mujeres en el ámbito comercial --un elemento sin precedente en ninguno de los acuerdos de Nueva Generación anteriores, incluido el propio Mercosur--.


\subsubsection{El Acuerdo de Asociación UE-Mercosur (EMPA) y el Acuerdo Interino de Comercio (iTA)}
Tras veinticinco años de negociación, iniciada en 1999, la Unión y los cuatro Estados parte del Mercosur (Argentina, Brasil, Paraguay y Uruguay) cerraron políticamente las negociaciones el 6 de diciembre de 2024, en la Cumbre de Montevideo. 

El resultado se articuló, siguiendo el modelo escindido ya consolidado, en dos instrumentos paralelos: (i) el Acuerdo de Asociación UE-Mercosur (EMPA), que combina diálogo político, cooperación y comercio, que por su contenido mixto requiere ratificación de los veintisiete Estados miembros; y, (ii) el Acuerdo Interino de Comercio (iTA), limitado a las materias de competencia exclusiva de la Unión, concebido para aplicarse antes de la entrada en vigor definitiva del EMPA y llamado a quedar derogado y sustituido por este una vez completada su ratificación. El Consejo autorizó la firma de ambos instrumentos el 9 de enero de 2026, y la firma tuvo lugar en Asunción el 17 de enero de 2026; el iTA se aplica provisionalmente desde el 1 de mayo de 2026, tras la correspondiente notificación al depositario de los tratados del Mercosur.\footnote{%
    Consejo de la UE, "Cómo funcionan los acuerdos UE-Mercosur", consilium.europa.eu; Comisión Europea, "Acuerdo de Asociación UE-Mercosur y Acuerdo Comercial interino", Access2Markets, 2026.
}

El acuerdo prevé la eliminación arancelaria sobre el $91\%$ de las exportaciones europeas al Mercosur en un plazo de hasta quince años (incluyendo el automóvil, gravado hasta ahora con un arancel del $35\%$, maquinaria y bebidas alcohólicas), mientras que la Unión abre progresivamente su mercado a productos agrícolas sensibles mediante el incremento de contingentes arancelarios (entre ellos 99.000 toneladas adicionales de carne de vacuno, además de aumentos en aves de corral, arroz, miel, soja y bioetanol), sujetos a un régimen de vigilancia reforzada. Se protegen, además, $350$ indicaciones geográficas europeas frente a imitaciones en los mercados del Mercosur, y se abre por primera vez para estos países, que no son signatarios del GPA de la OMC, un acceso relevante, aunque no equiparable al de los acuerdos con Canadá o Japón, a la contratación pública. Como contrapartida a la apertura del mercado agrícola europeo, el Consejo adoptó formalmente, el 5 de marzo de 2026, un Reglamento que activa cláusulas bilaterales de salvaguardia específicamente diseñadas para los productos agrícolas sensibles del acuerdo (vacuno, aves, porcino, azúcar, etanol, arroz, miel, maíz), que permite a la Comisión, a petición de un Estado miembro, iniciar investigaciones y adoptar medidas correctoras de manera acelerada si se produce un aumento sustancial de las importaciones o un desplome de las cotizaciones internas.\footnote{%
    Consejo de la UE, comunicado de prensa de 5 de marzo de 2026, "UE-Mercosur: el Consejo da luz verde a las salvaguardias para los productos agrícolas".
} Este mecanismo se ha activado ya, en la práctica, apenas semanas después de la aplicación provisional del acuerdo, ante la caída de las cotizaciones del vacuno en el mercado español.

A diferencia de los acuerdos con Canadá, Singapur, Vietnam o México, el acuerdo con Mercosur no incorpora un capítulo de protección de inversiones con mecanismo ICS, dada la resistencia histórica de los países del Mercosur al arbitraje inversor-Estado. En cambio, el capítulo de Comercio y Desarrollo Sostenible incorpora una innovación sin precedentes en la práctica convencional de la Unión: la designación del Acuerdo de París como elemento esencial del acuerdo global, en un régimen análogo al de la clásica cláusula democrática de derechos humanos, lo que permite la suspensión parcial o total del acuerdo si una de las partes se retira del Acuerdo de París o deja de aplicarlo de buena fe; se incorpora asimismo un compromiso jurídicamente vinculante de detener la deforestación a partir de 2030.\footnote{%
    Comisión Europea, Access2Markets, "Acuerdo de Asociación UE-Mercosur y Acuerdo Comercial interino"; Ministerio de Economía, Comercio y Empresa, "El Acuerdo UE-MERCOSUR".
} La efectividad real de esta cláusula es, no obstante, objeto de debate doctrinal: se señala que los mecanismos de suspensión de los capítulos TSD anteriores rara vez se han activado en la práctica, dado el elevado coste político y económico de recurrir a ellos.

\paragraph{Estado actual: Bloqueo procesal ante el TJUE} 
El 21 de enero de 2026, apenas cuatro días después de la firma, el Parlamento Europeo decidió remitir el acuerdo al TJUE para que se pronuncie sobre su compatibilidad con los Tratados, lo que ha paralizado de facto el procedimiento de ratificación del EMPA --sin afectar, no obstante, a la aplicación provisional del iTA, que descansa en su condición de acuerdo de competencia exclusiva--. Este episodio constituye el ejemplo más reciente y más citado del problema de los acuerdos mixtos ya expuesto en el epígrafe anterior.


\subsection{Acuerdos de segunda línea}
Los acuerdos que se agrupan bajo esta categoría comparten con los de «Nueva Generación» su naturaleza de acuerdos de libre comercio recíprocos, pero se distinguen de ellos por una intensidad de contenido sensiblemente menor: carecen, por lo general, de un capítulo de servicios con enfoque de lista negativa, no incorporan --o lo hacen de forma incompleta-- un mecanismo de solución de controversias en materia de inversiones del tipo ICS, y su capítulo TSD suele limitarse a compromisos de cooperación y de mejores esfuerzos, sin el aparato sancionador que progresivamente se ha ido incorporando a los acuerdos más recientes. Esta menor intensidad no obedece necesariamente a una menor ambición política de la Unión, sino, en muchos casos, al momento histórico en que se negociaron y al menor peso económico relativo de la contraparte --si bien, como demuestra el propio proceso de modernización mexicano, algunos de estos acuerdos están en trayectoria de converger, con el tiempo, hacia el modelo de Nueva Generación.

\subsubsection{Acuerdos Comerciales con Perú, Colombia y Ecuador}
El Acuerdo Comercial multipartito entre la Unión y Perú y Colombia, firmado en 2012, constituye --junto con el posterior acuerdo de adhesión de Ecuador en 2017-- el ejemplo más representativo de esta categoría. A diferencia de los Acuerdos de Asociación con Centroamérica o con los países andinos originalmente negociados como bloque regional, este acuerdo se fundamenta exclusivamente en el artículo 207 TFUE --sin componente de asociación política ni cooperación al desarrollo--, lo que explica su naturaleza estrictamente comercial y su tramitación más ágil.

El acuerdo elimina aranceles sobre la práctica totalidad del comercio de mercancías y liberaliza de forma significativa el comercio de servicios, pero mantiene un enfoque de lista positiva --y no negativa, como el CETA--, y no incorpora mecanismo ICS alguno para la protección de inversiones, que se rige en su lugar por los APPRI bilaterales preexistentes entre cada Estado miembro y cada país andino. El capítulo TSD, de redacción anterior al modelo reforzado de 2022, ha sido objeto de escrutinio específico en relación con Colombia, donde organizaciones sindicales han presentado reiteradas denuncias por violencia contra sindicalistas y por el incumplimiento de compromisos en materia de libertad sindical, sin que hasta la fecha se haya activado un panel de expertos equivalente al del caso coreano.

\subsubsection{Acuerdo de Asociación Económica, Concertación Política y Cooperación con México (Acuerdo Global, 2000)}
Las relaciones comerciales entre la Unión y México se rigen históricamente por el Acuerdo de Asociación Económica, Concertación Política y Cooperación el Acuerdo Global, en vigor desde el año 2000 y pionero, en su momento, como primer acuerdo de este tipo celebrado por la Unión con un país de América Latina.

Tras negociaciones iniciadas en 2016, la Unión y México alcanzaron la conclusión política de la modernización del Acuerdo Global en enero de 2025 y adoptaron el modelo ya consolidado de escisión en dos instrumentos: el Acuerdo Global Modernizado (AGM), de naturaleza mixta por incluir un pilar político y de cooperación además del comercial y de inversión, y un Acuerdo Comercial Interino (ACI), de competencia exclusiva de la Unión. Ambos se firmaron el 22 de mayo de 2026, en el marco de la VIII Cumbre México-Unión Europea, y el Parlamento Europeo prestó su aprobación a ambos --479 votos a favor para el AGM y 474 para el ACI-- a comienzos de julio de 2026.\footnote{%
    Parlamento Europeo, aprobación del Acuerdo Global Modernizado y del Acuerdo Comercial Interino con México, julio de 2026; Secretaría de Relaciones Exteriores de México, comunicado de 22 de mayo de 2026.
} El ACI podrá entrar en vigor tan pronto como concluya el proceso de ratificación por parte de México, sin esperar a la ratificación de los veintisiete Estados miembros que sí exige el AGM en su conjunto.

La modernización introduce, frente al Acuerdo Global de 2000, un capítulo de protección de inversiones, la apertura recíproca de la contratación pública, un capítulo de comercio digital, la protección de un número relevante de indicaciones geográficas europeas y un capítulo TSD reforzado conforme al nuevo enfoque de 2022. A diferencia de Chile, sin embargo, el AGM mexicano no incorpora un capítulo de género y comercio ni alcanza el mismo grado de apertura arancelaria, lo que --pese a su contenido reforzado-- justifica mantenerlo, por ahora, en esta categoría intermedia. Debe notarse, además, que el propio Parlamento Europeo ha condicionado políticamente su respaldo a la vigilancia de la situación de violencia política y de derechos humanos en México, lo que se ha traducido en una resolución específica de acompañamiento aprobada junto con el propio Acuerdo.\footnote{%
    Resolución del Parlamento Europeo A-10-2026-0182, sobre la situación de derechos humanos y violencia política en México en el contexto del Acuerdo Global Modernizado.
}

\subsubsection{Acuerdo Euromediterráneo de Asociación con Israel}
El Acuerdo Euromediterráneo de Asociación con Israel, firmado en 1995 en el marco del Proceso de Barcelona y en vigor desde el año 2000, liberaliza el comercio de productos industriales y contiene disposiciones limitadas sobre productos agrícolas. Como en todos los Acuerdos de Asociación de la Unión, su artículo 2 configura el respeto de los derechos humanos y de los principios democráticos como «elemento esencial» del Acuerdo (la misma técnica jurídica, análoga a una cláusula democrática, que hemos visto incorporarse por primera vez a un ámbito distinto de los derechos humanos en el capítulo climático del acuerdo con Mercosur).

\paragraph{Revisión del cumplimiento (\textit{ex} art. 2: Conflicto en Gaza} 
A raíz del conflicto en la Franja de Gaza, el Alto Representante de la Unión presentó al Consejo de Asuntos Exteriores, el 23 de junio de 2025, una revisión que concluyó la existencia de indicios de incumplimiento por parte de Israel del artículo 2 del Acuerdo.\footnote{%
    Consejo de Asuntos Exteriores de la UE, 23 de junio de 2025; Centro de Documentación Europea, Universidad de Granada.
} En septiembre de 2025, la presidenta de la Comisión anunció una propuesta de suspensión parcial de las disposiciones comerciales del Acuerdo, junto con sanciones específicas dirigidas a determinados ministros y a colonos vinculados a episodios de violencia en Cisjordania. A la fecha de las últimas informaciones disponibles (mayo de 2026), esa propuesta seguía sin obtener la mayoría necesaria en el Consejo, bloqueada singularmente por Alemania e Italia, por lo que el Acuerdo permanece formalmente en vigor en su integridad.\footnote{%
    Amnistía Internacional, "Líneas rojas, no alfombras rojas", mayo de 2026; se trata de un asunto políticamente controvertido y en evolución activa, sobre el que distintos actores institucionales y organizaciones mantienen posiciones encontradas.
} Este caso constituye, junto con el más antiguo precedente de la jurisprudencia \textit{Brita} (2010) --sobre el origen preferencial de mercancías procedentes de los asentamientos-- y la sentencia \textit{Psagot} (2019) --sobre el etiquetado de origen de productos de los asentamientos--, uno de los ejemplos más citados en la doctrina para ilustrar la tensión entre la cláusula de derechos humanos de los Acuerdos de Asociación y su aplicación práctica efectiva.

\subsection{Acuerdos de Asociación Económica (EPA): África, Caribe y Pacífico (ACP)}
Los Acuerdos de Asociación Económica (EPA, por sus siglas en inglés) constituyen la respuesta de la Unión a un problema jurídico concreto planteado en el seno de la OMC. 

Las preferencias comerciales no recíprocas que la Unión concedía a los países de África, el Caribe y el Pacífico bajo los sucesivos Convenios de Yaundé (1963-1975) y Lomé (1975-2000) eran, en principio, incompatibles con la cláusula de nación más favorecida del artículo I del GATT, al beneficiar selectivamente a un grupo de países en desarrollo y no a la totalidad de estos, lo que impedía ampararlas en la cláusula de habilitación (\textit{enabling clause}), aplicable únicamente a esquemas de preferencias universales como el SPG. Su mantenimiento solo fue posible gracias a sucesivas exenciones (\textit{waivers}) del artículo I del GATT concedidas por la OMC, la última de ellas obtenida en la Conferencia Ministerial de Doha de 2001, con vencimiento el 31 de diciembre de 2007.\footnote{%
    Sobre la incompatibilidad de las preferencias de Lomé con el GATT y el mecanismo del \textit{waiver}, véase FAO, "La situación de las preferencias comerciales en la OMC"; y el litigio \textit{Comunidades Europeas -- Régimen para la importación, venta y distribución de bananos} (DS27), donde el Órgano de Apelación de la OMC constató la pérdida de cobertura de la exención de Doha aplicada a los bananos ACP a partir del 1 de enero de 2006.
} 

El Acuerdo de Cotonú (2000), que sucedió a Lomé, prorrogó las preferencias no recíprocas únicamente durante un período preparatorio (2000-2007), con el mandato expreso de sustituirlas, antes de esa fecha, por acuerdos comerciales recíprocos y compatibles con las normas de la OMC: los EPA.

\paragraph{El principio de asimetría.} 
A diferencia de los acuerdos de «Nueva Generación» y de «segunda línea», la reciprocidad exigida por la OMC en los EPA se articula de forma marcadamente asimétrica a favor del socio ACP: la Unión abre su mercado de forma inmediata, plena y prácticamente sin excepciones, mientras que los países ACP disponen de periodos de transición muy dilatados (entre quince y veinticinco años) para liberalizar sus importaciones procedentes de la Unión, pueden excluir productos sensibles de la liberalización, y se benefician de cláusulas de salvaguardia especiales vinculadas a la seguridad alimentaria y a la protección de industrias incipientes.

\paragraph{Nuevo Marco Político: Acuerdo de Samoa (2023)} 
El marco político general que envuelve a los EPA ha cambiado recientemente: el Acuerdo de Cotonú, agotado su período de vigencia, ha sido sustituido por el Acuerdo de Samoa, firmado el 15 de noviembre de 2023 entre la Unión y los $79$ miembros de la Organización de Estados de África, el Caribe y el Pacífico (OEACP), y aplicado provisionalmente desde el 1 de enero de 2024.

El Acuerdo de Samoa constituye un marco común aplicable a los $79$ Estados, complementado por tres protocolos regionales específicos para África, el Caribe y el Pacífico, y regirá la relación política, de cooperación y de derechos humanos durante los próximos veinte años; no sustituye, sin embargo, a los EPA regionales, que permanecen en vigor como pilar estrictamente comercial de la relación, ahora enmarcado bajo el paraguas político del nuevo Acuerdo.

\subsubsection{El AAE CARIFORUM-UE: Modelo de referencia}
El Acuerdo de Asociación Económica (2008) con los catorce Estados del CARIFORUM (los miembros de la Comunidad del Caribe más República Dominicana), fue el primer EPA celebrado por la Unión, y se ha convertido en el modelo de referencia para las negociaciones posteriores con las demás regiones ACP.

A diferencia del resto de EPA, que en su mayoría se han limitado durante años al comercio de mercancías, este incorpora desde su origen compromisos sustanciales en materia de comercio de servicios (con los Estados más desarrollados del CARIFORUM comprometiéndose a liberalizar el $75\%$ de su comercio de servicios con la Unión, y los menos desarrollados el $65\%$, inversiones, política de competencia, contratación pública y propiedad intelectual, además de un Protocolo específico de cooperación cultural. Pese a esta profundidad de contenido, su fundamento último sigue siendo asimétrico: el acceso de la Unión es inmediato, mientras que los Estados del CARIFORUM disponen de hasta veinticinco años para liberalizar sus importaciones procedentes de la Unión.

\subsubsection{Comunidad para el Desarrollo del África Meridional (SADC, 2016)}
El AAE UE-SADC, firmado el 10 de junio de 2016 por Botsuana, Esuatini, Lesoto, Mozambique, Namibia y Sudáfrica, se aplica provisionalmente desde el 10 de octubre de 2016.

\paragraph{Caso particular: Sudáfrica} 
Este acuerdo reviste una importancia particular por haber sustituido, en su vertiente estrictamente comercial, al Acuerdo de Comercio, Desarrollo y Cooperación (TDCA) que la Unión mantenía bilateralmente con Sudáfrica desde 1999. 

Sudáfrica ocupa dentro del EPA de la SADC una posición singular y menos favorable que el resto de socios. Por un lado, mientras que Botsuana, Esuatini, Lesoto, Mozambique y Namibia (de la Unión Aduanera del África Austral, SACU) reciben de la Unión un acceso libre de aranceles y contingentes para prácticamente la totalidad de sus exportaciones (salvo armas y municiones), Sudáfrica, por su mayor nivel de desarrollo relativo, solo obtiene ese trato para el $96\%$ de sus exportaciones, beneficiándose del restante únicamente de una reducción arancelaria o de contingentes preferentes. En sentido inverso, el conjunto SACU abre su mercado a la Unión de forma libre de aranceles y contingentes solo para el $84,9\%$ de las importaciones europeas. 

El acuerdo prevé además la posibilidad de que las partes lo profundicen en el futuro hacia servicios, inversiones, comercio digital y comercio y desarrollo sostenible.

\subsubsection{Los EPA africanos de implementación fragmentada}
Frente al caso relativamente ordenado del CARIFORUM y de la SADC, el resto de configuraciones regionales africanas del EPA ilustran, más que un régimen consolidado, un proceso de implementación fragmentado y, en buena parte, todavía inconcluso.

Por un lado, el AAE con los dieciséis Estados de África Occidental (CEDEAO) negociado como bloque regional, no se aplica en la actualidad en ninguno de los países de la región. La razón es que su entrada en aplicación provisional exige que la totalidad de los Estados lo hayan firmado y que al menos dos tercios lo hayan ratificado, condición que continúa sin cumplirse porque Nigeria no lo ha firmado, alegando el riesgo de desindustrialización que supondría para su sector manufacturero la apertura recíproca frente a la Unión.

En África Central únicamente Camerún aplica, desde 2014, un EPA interino y limitado al comercio de mercancías, mientras que el resto de Estados de la región permanecen fuera del acuerdo y se benefician, en su lugar, del régimen unilateral del SPG o, en el caso de los países menos desarrollados, del régimen EBA.

Las negociaciones del AAE con la Comunidad del África Oriental (EAC) concluyeron en octubre de 2014, y el acuerdo ha sido firmado por Kenia, por Ruanda y por todos los Estados miembros de la Unión, pero continúa sin entrar en vigor a la espera de la firma del resto de socios de la EAC. Este bloqueo resulta especialmente ilustrativo porque, al tratarse de un acuerdo con una unión aduanera regional ya constituida (el Arancel Exterior Común de la EAC), su falta de entrada en vigor genera además fricciones internas dentro del propio bloque regional entre los Estados que sí querrían aplicarlo --Kenia, principal exportador de la región a la Unión y único Estado de la EAC no clasificado como país menos desarrollado, por lo que sin el EPA perdería su acceso preferente al mercado europeo-- y los que se resisten a hacerlo.

\paragraph{Crítica: Implementación asimétrica}
Esta implementación desigual es señalada de forma recurrente por la doctrina como evidencia de que el formato de negociación por bloques regionales, pensado para reforzar la integración económica sur-sur, ha topado en la práctica con las asimetrías de intereses económicos entre los propios Estados de cada bloque, retrasando los beneficios que el propio EPA pretendía generar para el conjunto de la región.

\subsubsection{AAE con los Estados del Pacífico}
El AAE interino con los Estados ACP del Pacífico fue firmado inicialmente por Papúa Nueva Guinea (julio de 2009) y por Fiyi (diciembre de 2009), aplicándose provisionalmente Papúa Nueva Guinea tras su ratificación en 2011 y Fiyi desde 2014; posteriormente se han adherido Samoa (diciembre de 2018) y las Islas Salomón (mayo de 2020). De los catorce Estados ACP de la región, Papúa Nueva Guinea y Fiyi concentran, con diferencia, el grueso del comercio bilateral con la Unión, lo que explica que el resto de Estados insulares del Pacífico --de reducidísimo peso comercial-- no hayan mostrado el mismo interés en completar su adhesión al Acuerdo.


\subsection{Sistema de Preferencias Generalizadas (SPG)}
El SPG es, de entre todos los regímenes estudiados en esta sección, el único que carece de fundamento convencional: no nace de un tratado internacional negociado con un tercer país, sino de un Reglamento de la Unión --acto unilateral autónomo, adoptado sobre la base del artículo 207 TFUE-- mediante el que la Unión concede, por propia iniciativa y sin obligación de reciprocidad, una reducción o supresión de sus derechos de aduana a los productos originarios de países en desarrollo. Su compatibilidad con las normas de la OMC se sustenta, a diferencia de las preferencias históricas ACP ya estudiadas, en la cláusula de habilitación (\textit{enabling clause}) de 1979, que sí ampara --a diferencia de lo que ocurría con Lomé-- los esquemas de preferencias arancelarias generalizadas, no discriminatorios entre países en desarrollo, y sin necesidad de exención específica alguna.

\paragraph{Nuevo Reglamento: Reglamento (UE) 2026/1395} 
El régimen ha estado gobernado, desde 2014, por el Reglamento (UE) n.º 978/2012, cuyo periodo de aplicación fue objeto de sucesivas prórrogas técnicas, la última de ellas hasta diciembre de 2027, mientras se tramitaba su reforma. 

El nuevo Reglamento (UE) 2026/1395, aprobado por el Parlamento Europeo el 28 de abril de 2026 y publicado en el Diario Oficial el 22 de junio de 2026, sustituye al Reglamento de 2012 y será aplicable desde el 1 de enero de 2027, por un nuevo periodo de diez años. El sistema mantiene, en su estructura básica, los tres regímenes ya conocidos, articulados en una jerarquía ascendente de intensidad de la preferencia.

\subsubsection{Régimen General: SPG}
El régimen general del SPG se dirige a todos los países en desarrollo que, no beneficiándose de ningún régimen más favorable, comparten una necesidad de desarrollo común y se encuentran en una fase similar de desarrollo económico, y consiste en la reducción de los derechos de aduana sobre productos industriales y agrícolas no sensibles.

A diferencia del régimen EBA, el SPG estándar incorpora un mecanismo de \textit{graduación} que excluye automáticamente de sus beneficios determinados países y productos. En primer lugar, a aquellos países que el Banco Mundial clasifica como de renta alta o media-alta durante tres años consecutivos.\footnote{%
    Criterio en cuya aplicación, por ejemplo, Indonesia perderá su condición de beneficiario del SPG a partir del 1 de enero de 2027, tras haber sido clasificada como país de renta media-alta en 2023, 2024 y 2025.
} Y, en segundo lugar, a las secciones de productos en que un país beneficiario alcanza una cuota de mercado dominante en las importaciones de la Unión bajo el SPG durante tres años consecutivos, lo que da lugar a la suspensión sectorial de las preferencias correspondientes (no del país en su conjunto).

\subsubsection{Régimen especial: Desarrollo Sostenible y la Gobernanza (SPG+)}
El SPG+ concede, a los países que se acogen a él, la eliminación completa de los derechos de aduana sobre prácticamente las mismas partidas que cubre el régimen general, a cambio de la ratificación y la aplicación efectiva de un listado cerrado de convenios internacionales en materia de derechos humanos, derechos laborales, medio ambiente y buena gobernanza --entre ellos la Declaración de las Naciones Unidas sobre el Derecho al Desarrollo (1986), la Declaración de Río sobre el Medio Ambiente y el Desarrollo (1992), la Declaración de la OIT relativa a los Principios y Derechos Fundamentales en el Trabajo (1998), la Declaración del Milenio de las Naciones Unidas (2000) y la Declaración de Johannesburgo sobre el Desarrollo Sostenible (2002)--.

\paragraph{Un régimen sujeto a condicionalidad y monitorización continuada.} A diferencia del régimen general, el acceso al SPG+ no es automático, sino que exige solicitud expresa del país interesado y una evaluación previa de vulnerabilidad económica --escasa diversificación exportadora e importaciones limitadas de la Unión--, y su mantenimiento está sujeto a un ciclo de monitorización bienal por parte de la Comisión sobre el grado de cumplimiento efectivo de los veintisiete convenios. Este mecanismo de condicionalidad no es meramente formal: en agosto de 2020 la Unión retiró parcialmente --por primera vez en la historia del sistema-- las preferencias EBA a Camboya, suspendiendo el acceso libre de aranceles para determinadas partidas de prendas de vestir, calzado y productos de viaje, ante violaciones graves y sistemáticas de derechos civiles y políticos y de la libertad sindical, lo que ilustra que la condicionalidad de estos regímenes especiales, aunque de aplicación poco frecuente, no es puramente retórica.

\subsubsection{El régimen especial a favor de los países menos desarrollados: «Todo menos armas» (EBA)}

El régimen EBA concede a los países clasificados por Naciones Unidas como países menos adelantados (PMA) un acceso libre de aranceles y de contingentes al mercado de la Unión para la totalidad de sus productos, sin excepción, salvo las armas y municiones. A diferencia del SPG estándar, el EBA carece de mecanismo de graduación por producto y de límite temporal de aplicación, y --a diferencia también del SPG estándar-- un país no pierde automáticamente su condición de beneficiario EBA por el hecho de celebrar un acuerdo de libre comercio con la Unión; su inclusión y exclusión depende exclusivamente de la clasificación de Naciones Unidas como PMA, sin necesidad de solicitud alguna por parte del país beneficiario.

\paragraph{El caso de Bangladesh: la prueba de fuego del sistema.} Bangladesh --que concentra, con diferencia, la mayor cuota de beneficio del régimen EBA a nivel mundial, siendo el origen de aproximadamente dos tercios de las importaciones globales acogidas a este régimen-- está previsto que se gradúe formalmente de la categoría de PMA de Naciones Unidas el 24 de noviembre de 2026.\footnote{%
    Naciones Unidas, Portal de los Países Menos Adelantados, "Bangladesh graduation status".
} Conforme a la práctica habitual de la Unión, la graduación no implica la pérdida inmediata de las preferencias EBA, sino la apertura de un período de transición de tres años --hasta noviembre de 2029--, durante el cual Bangladesh seguirá beneficiándose del régimen; una vez concluido ese plazo, y a falta de negociación de un acuerdo bilateral, sus exportaciones quedarán sujetas al régimen general del SPG --con tipos arancelarios en torno al 9,6\,\% para buena parte de sus exportaciones textiles, principal partida exportadora del país-- salvo que logre acreditar el cumplimiento de los veintisiete convenios exigidos para acceder al SPG+.\footnote{%
    The Daily Star, "EU's GSP+: The lifeline Bangladesh must win before 2029", 2026; estimaciones de la OMC citadas en dicha fuente sitúan en torno a 8.000 millones de dólares anuales --cerca del 14\,\% de las exportaciones totales del país-- el impacto potencial de la pérdida de las preferencias EBA.
} Este caso constituye, en la actualidad, el ejemplo más citado para ilustrar la lógica de graduación progresiva --y sus riesgos económicos asociados-- que estructura todo el sistema del SPG.

\paragraph{La erosión estructural del sistema.} Como se ha apuntado ya al hilo de los demás regímenes preferenciales, la proliferación de acuerdos de libre comercio bilaterales de la Unión --de Nueva Generación, de segunda línea o EPA-- erosiona progresivamente el valor relativo de las preferencias unilaterales del SPG: a medida que más socios comerciales obtienen, vía acuerdo negociado, un acceso equivalente o superior al del SPG, el margen de preferencia --y, con él, el incentivo real que syrve al país beneficiario del SPG para no perseguir un acuerdo bilateral propio-- se reduce, en lo que la doctrina describe como el coste estructural, y en cierto modo paradójico, del propio éxito de la política de acuerdos comerciales de la Unión.


























\clearpage
\section{Conclusión} 


\subsection{Recapitulación}


\subsection{Extensiones y relación con otras partes del Temario}



\subsection{Opinión}

\subsubsection{Idea final (Salida o cierra)}

\clearpage
\pagenumbering{gobble}
\MyBibliography

\href{https://www.cuatrecasas.com/es/spain/competencia-derecho-ue/art/fsr-contratacion-publica-primera-decision-con-condiciones}{\textit{FSR: primera decisión con condiciones en contratación pública}. Cuatrecasas, 2026}

\href{https://infoeuropa.uv.es/directrices-reglamento-subvenciones-extranjeras/}{\textit{Directrices relativas a la aplicación de determinadas disposiciones del Reglamento 2022/2560}. InfoEUROPA, Centro de Documentación Europea de Valencia, 2026}

\href{https://www.uria.com/es/publicaciones/8688-el-regimen-de-autorizacion-de-subvenciones-extranjeras-entra-en-aplicacion}{\textit{El régimen de autorización de subvenciones extranjeras entra en aplicación}. Uría Menéndez, 2023}

\href{https://eur-lex.europa.eu/legal-content/ES/ALL/?uri=CELEX%3A52026DC0368}{\textit{Informe sobre el Reglamento sobre las subvenciones extranjeras}. EUR-Lex (Comisión Europea), 2026}

\href{https://commission.europa.eu/strategy-and-policy/priorities-2019-2024/europe-fit-digital-age/european-industrial-strategy/foreign-subsidies-regulation_es}{\textit{Reglamento sobre subvenciones extranjeras}. Comisión Europea, 2026}

\href{https://www.euskadi.eus/gobierno-vasco/-/nota_prensa/2026/politica-de-competencia-la-ue-ha-clarificado-procedimientos-sobre-la-regulacion-de-subvenciones-extranjeras-que-afecten-a-concursos-publicos-o-concentraciones-monopolisticas-y-ha-validado-diversas-agregaciones-empresariales-en-ultimas-fechas/}{\textit{Política de competencia: la UE ha clarificado procedimientos sobre la regulación de subvenciones extranjeras}. Gobierno Vasco, 2026}

\href{https://periscopiofiscalylegal.pwc.es/publicacion-del-reglamento-europeo-sobre-subvenciones-extranjeras-que-distorsionan-el-mercado-interior/}{\textit{Publicación del Reglamento Europeo sobre Subvenciones Extranjeras que Distorsionan el Mercado Interior}. PwC Periscopio Fiscal y Legal, 2025}

\href{https://www.garrigues.com/es_ES/noticia/entra-vigor-reglamento-ue-subvenciones-extranjeras}{\textit{Entra en vigor el reglamento de la UE sobre subvenciones extranjeras}. Garrigues, 2023}

\href{https://www.boe.es/buscar/doc.php?id=DOUE-L-2022-81950}{\textit{Reglamento (UE) 2022/2560 del Parlamento Europeo y del Consejo, de 14 de diciembre de 2022}. BOE, 2022}

\href{https://noticias.juridicas.com/base_datos/Privado/744356-regl-2022-2560-ue-de-14-dic-subvenciones-extranjeras-que-distorsionan-el.html}{\textit{Reglamento (UE) 2022/2560, sobre las subvenciones extranjeras que distorsionan el mercado interior}. Noticias Jurídicas, 2022}

\href{https://www.europarl.europa.eu/meetdocs/2014_2019/plmrep/COMMITTEES/INTA/PR/2022/04-25/1245578ES.pdf}{\textit{Proyecto de informe sobre la Propuesta de Reglamento sobre las subvenciones extranjeras}. Christophe Hansen (Comisión de Comercio Internacional, Parlamento Europeo), 2021}

\href{https://sede.agenciatributaria.gob.es/Sede/aduanas/prohibiciones-restricciones-operaciones-comercio-exterior-carbono/mecanismo-ajuste-frontera-carbono.html}{\textit{Mecanismo de Ajuste en Frontera por Carbono (CBAM)}. Agencia Tributaria, 2026}

\href{https://www.orkestra.deusto.es/es/actualidad/noticias-eventos/beyondcompetitiveness/2752-mecanismo-de-ajuste-por-carbono-en-frontera-cbam}{\textit{¿Qué es el Mecanismo de Ajuste por Carbono en Frontera (CBAM) y cómo funciona?}. Orkestra Instituto Vasco de Competitividad, 2025}

\href{https://vatgreentax.com/precio-cbam-2026/}{\textit{Mecanismo de Ajuste en Frontera por Carbono (CBAM): todo lo que necesitas saber para la fase definitiva}. VAT Green Tax, 2026}

\href{https://www.ecogestor.com/cbam-mecanismo-ajuste-frontera-carbono/}{\textit{El CBAM entra en su fase definitiva: normativa consolidada y efectos clave para 2026}. Eurofins EcoGestor, 2026}

\href{https://www.preferredbynature.org/es/EUDR}{\textit{Reglamento EUDR}. Preferred by Nature, 2026}

\href{https://actualidad.aidimme.es/2025/12/18/la-ue-aprueba-un-nuevo-aplazamiento-del-reglamento-eudr-por-un-ano-alivia-la-gestion-y-otorga-a-las-pyme-6-meses-mas-de-gracia/}{\textit{La UE aprueba un nuevo aplazamiento del Reglamento EUDR por un año}. AIDIMME, 2025}

\href{https://www.asd-int.com/es/mini_blog/reglamento-deforestacion-eudr-entrada-en-vigor-aplazada/}{\textit{EUDR: aplazamiento de la entrada en vigor}. ASD International, 2026}

\href{https://www.miteco.gob.es/es/biodiversidad/temas/politica-forestal/madera-legal-productos-libres-defor/eudr/-que-es-el-eudr-/calendario-de-aplicacion.html}{\textit{Calendario de aplicación del EUDR}. Ministerio para la Transición Ecológica y el Reto Demográfico (MITECO), 2026}

\href{https://ejaso.com/conocimiento/reglamento-de-deforestacion-eudr-y-modificacion-del-regimen-cbam}{\textit{Reglamento de Deforestación (EUDR) y modificación del régimen CBAM}. EJASO, 2025}

\href{https://trade.ec.europa.eu/access-to-markets/es/news/la-aplicacion-del-reglamento-eudr-sobre-productos-libres-de-deforestacion-se-retrasa-hasta}{\textit{La aplicación del Reglamento EUDR sobre productos libres de deforestación se retrasa hasta diciembre de 2025}. Access2Markets, Comisión Europea, 2024}

\href{https://dialnet.unirioja.es/servlet/articulo?codigo=8996373}{\textit{El "arancel al carbono (CBAM)": ¿proteccionismo verde o liderazgo global contra el cambio climático?}. Dialnet, 2023}

\href{https://www.realinstitutoelcano.org/analisis/el-arancel-al-carbono-cbam-proteccionismo-verde-o-liderazgo-global-contra-el-cambio-climatico/}{\textit{El "arancel al carbono (CBAM)": ¿proteccionismo verde o liderazgo global contra el cambio climático?}. Enrique Feás, Federico Steinberg y Lara Lázaro Touza, Real Instituto Elcano, 2023}

\href{https://www.wto.org/spanish/tratop_s/envir_s/envt_rules_exceptions_s.htm}{\textit{Normas de la OMC y políticas ambientales: excepciones previstas en el GATT}. Organización Mundial del Comercio, s.f.}

\href{https://www.wto.org/spanish/tratop_s/dispu_s/repertory_s/g3_s.htm}{\textit{Repertorio del Órgano de Apelación de la OMC: Informes y Laudos (1995-2005)}. Organización Mundial del Comercio, 2005}

\href{https://www.eleconomista.es/economia/noticias/13905474/05/26/francia-pide-activar-el-instrumento-anticoercion-ante-la-nueva-amenaza-de-trump-de-aplicar-aranceles-del-20.html}{\textit{Francia pide activar el instrumento anticoerción ante la nueva amenaza de Trump}. El Economista, 2026}

\href{https://www.infobae.com/espana/agencias/2026/01/18/la-ue-mas-cerca-que-nunca-de-activar-el-instrumento-anticoercion-por-las-amenazas-de-trump/}{\textit{La UE más cerca que nunca de activar el instrumento anticoerción por las amenazas de Trump}. Infobae (EFECOM), 2026}

\href{https://donporque.com/que-es-el-instrumento-anticoercion/}{\textit{¿Qué es el Instrumento Anticoerción, la bazuca de la UE?}. Don Porqué, 2026}

\href{https://en.wikipedia.org/wiki/Anti-Coercion_Instrument}{\textit{Anti-Coercion Instrument}. Wikipedia, 2026}

\href{https://policy.trade.ec.europa.eu/enforcement-and-protection/protecting-against-coercion_en}{\textit{Protecting against coercion}. Trade and Economic Security, Comisión Europea, 2026}

\href{https://es.euronews.com/2026/03/06/sanchez-la-ue-se-arman-frente-trump-asi-es-el-articulo-5-mecanismo-secreto-de-la-ue}{\textit{Así es el instrumento anti-coerción de la UE, el arma disuasoria que Sánchez tendría frente a Trump}. Euronews, 2026}

\href{https://www.swissinfo.ch/spa/claves-del-instrumento-de-la-ue-contra-la-coerci%C3%B3n-econ%C3%B3mica-que-podr%C3%ADa-usar-contra-ee.uu./89105365}{\textit{Claves del instrumento de la UE contra la coerción económica que podría usar contra EE.UU.}. SWI swissinfo.ch (EFE), 2025}

\href{https://www.cidob.org/publicaciones/el-instrumento-anticoercion-propuesto-por-la-comision-europea-desde-el-punto-de-vista}{\textit{El instrumento anticoerción propuesto por la Comisión Europea desde el punto de vista del derecho internacional}. CIDOB, s.f.}

\href{https://trade.ec.europa.eu/access-to-markets/es/content/instrumento-contra-la-coercion}{\textit{Instrumento contra la coerción}. Access2Markets, Comisión Europea, 2026}

\href{https://e-archivo.uc3m.es/bitstreams/aedef20b-ef53-4bb2-93b1-128d15089a5e/download}{\textit{El Reglamento sobre obstáculos al comercio}. e-Archivo, Universidad Carlos III de Madrid, 2015}

\href{https://slideplayer.es/slide/4152497/}{\textit{El Reglamento Obstáculos al Comercio (ROC)}. SlidePlayer, s.f.}

\href{https://comercio.gob.es/Barreras_Comerciales/barreras-terceros-paises/Paginas/ROC.aspx}{\textit{Reglamento de obstáculos al comercio (ROC)}. Ministerio de Economía, Comercio y Empresa, 2026}

\href{https://eur-lex.europa.eu/legal-content/ES/TXT/?uri=LEGISSUM:r11007}{\textit{Defensa contra los obstáculos al comercio}. EUR-Lex, 2014}

\href{https://www.consilium.europa.eu/es/policies/eu-mercosur-agreements-explained/}{\textit{Cómo funcionan los acuerdos UE-Mercosur}. Consejo de la Unión Europea, 2026}

\href{https://www.consilium.europa.eu/es/infographics/eu-mercosur-trade/}{\textit{Comercio UE-Mercosur: datos y cifras}. Consejo de la Unión Europea, 2026}

\href{https://www.consilium.europa.eu/es/press/press-releases/2026/01/09/eu-mercosur-council-greenlights-signature-of-the-comprehensive-partnership-and-trade-agreement/}{\textit{UE-Mercosur: el Consejo da luz verde a la firma del acuerdo de asociación y de comercio}. Consejo de la Unión Europea, 2026}

\href{https://commission.europa.eu/topics/trade/eu-mercosur-trade-agreement_es}{\textit{El acuerdo comercial UE-Mercosur}. Comisión Europea, 2026}

\href{https://es.wikipedia.org/wiki/Acuerdo_comercial_Mercosur-Uni%C3%B3n_Europea}{\textit{Acuerdo comercial Mercosur-Unión Europea}. Wikipedia, 2026}

\href{https://www.cope.es/emisoras/la-rioja/la-rioja-provincia/logrono/noticias/son-clausulas-salvaguarda-ue-refuerza-defensas-agricolas-acuerdo-mercosur-activara-medidas-suben-importaciones-caen-precios-20260211_3305384.amp.html}{\textit{¿Qué son las cláusulas de salvaguarda? La UE refuerza las defensas agrícolas ante el acuerdo con Mercosur}. COPE, 2026}

\href{https://agroinformacion.com/plantean-la-activacion-de-la-clausula-de-salvaguardia-contemplada-en-el-acuerdo-ue-mercosur-por-el-desplome-de-las-cotizaciones-del-vacuno/}{\textit{Plantean la activación de la cláusula de salvaguardia contemplada en el Acuerdo UE-Mercosur}. Agroinformación, 2026}

\href{https://www.europarl.europa.eu/news/es/press-room/20251211IPR32162/mercosur-el-parlamento-aprueba-medidas-de-salvaguarda-para-productos-agricolas}{\textit{Mercosur: el Parlamento aprueba medidas de salvaguarda para productos agrícolas}. Parlamento Europeo, 2025}

\href{https://trade.ec.europa.eu/access-to-markets/es/content/acuerdo-de-asociacion-ue-mercosur-y-acuerdo-comercial-interino}{\textit{Acuerdo de Asociación UE-Mercosur y Acuerdo Comercial interino}. Access2Markets, Comisión Europea, 2026}

\href{https://comercio.gob.es/PoliticaComercialUE/AcuerdosComerciales/Paginas/Mercosur.aspx}{\textit{El Acuerdo UE–MERCOSUR}. Ministerio de Economía, Comercio y Empresa, 2026}

\href{https://es.euronews.com/my-europe/2026/01/09/la-ue-da-luz-verde-al-acuerdo-con-el-mercosur-pese-al-rechazo-de-francia}{\textit{La UE da luz verde al acuerdo con el Mercosur pese al rechazo de Francia}. Euronews, 2026}

\href{https://www.eeas.europa.eu/delegations/argentina/acuerdo-ue-%E2%80%93-mercosur_und_en}{\textit{Acuerdo UE – Mercosur}. Servicio Europeo de Acción Exterior (EEAS), 2026}

\href{https://nuso.org/articulo/311-UE-y-mercosur-cuatro-nudos-un-desenlace/}{\textit{Unión Europea y Mercosur: cuatro nudos ¿y un desenlace?}. Nueva Sociedad, 2024}

\href{https://www.democrata.es/internacional/turquia-y-la-ue-dan-pasos-para-modernizar-la-union-aduanera-y-que-alcance-todo-su-potencial/}{\textit{Turquía y UE aceleran la modernización de la Unión Aduanera}. El Demócrata, 2026}

\href{https://www.infobae.com/america/agencias/2026/02/06/turquia-pide-actualizar-la-union-aduanera-con-la-ue-ante-la-incertidumbre-geopolitica/}{\textit{Turquía pide actualizar la unión aduanera con la UE ante la "incertidumbre geopolítica"}. Infobae (EFE), 2026}

\href{https://eur-lex.europa.eu/ES/legal-content/summary/eu-customs-union-with-turkey.html}{\textit{Unión Aduanera de la Unión Europea con Turquía}. EUR-Lex, 2016}

\href{https://www.infobae.com/america/agencias/2026/06/05/la-ue-busca-en-montenegro-un-nuevo-impulso-para-relanzar-la-ampliacion-a-los-balcanes/}{\textit{La UE busca en Montenegro un nuevo impulso para relanzar la ampliación a los Balcanes}. Infobae, 2026}

\href{https://balk.hu/es/2025/12/05/montenegro-albania-moldova-ukrajna-eu-bovites/}{\textit{La UE podría admitir a Montenegro, Albania, Moldavia y Ucrania lo antes posible}. Revista BALK, 2025}

\href{https://es.wikipedia.org/wiki/Ampliaci%C3%B3n_potencial_de_la_Uni%C3%B3n_Europea}{\textit{Ampliación potencial de la Unión Europea}. Wikipedia, 2026}

\href{https://www.moncloa.com/2026/06/04/ampliacion-ue-balcanes-cumbre-montenegro-3382141}{\textit{La UE acelera la ampliación hacia los Balcanes en una cumbre clave en Montenegro}. Moncloa.com, 2026}

\href{https://www.cidob.org/publicaciones/la-ampliacion-de-la-ue-los-balcanes-occidentales-urgencia-y-reestructuracion}{\textit{La ampliación de la UE a los Balcanes occidentales: urgencia y reestructuración}. CIDOB, s.f.}

\href{https://european-union.europa.eu/principles-countries-history/eu-enlargement_es}{\textit{Ampliación de la UE}. Unión Europea, 2026}

\href{https://es.wikipedia.org/wiki/Estrategia_para_los_Balcanes_Occidentales}{\textit{Estrategia para los Balcanes Occidentales}. Wikipedia, 2026}

\href{https://www.industriall-union.org/es/campana-sindical-lleva-a-corea-del-sur-a-ratificar-los-convenios-fundamentales-de-la-oit}{\textit{Campaña sindical lleva a Corea del Sur a ratificar los convenios fundamentales de la OIT}. IndustriALL Global Union, 2021}

\href{https://www.eleconomista.es/economia/noticias/13404061/06/25/el-acuerdo-marco-avanzado-entre-la-ue-y-chile-entra-en-vigor-de-manera-provisional.html}{\textit{El Acuerdo Marco Avanzado entre la UE y Chile entra en vigor de manera provisional}. El Economista (EFE), 2025}

\href{https://eur-lex.europa.eu/ES/legal-content/summary/eu-chile-advanced-framework-agreement.html?fromSummary=07}{\textit{Acuerdo Marco Avanzado Unión Europea-Chile}. EUR-Lex, 2025}

\href{https://www.ceoe.es/sites/ceoe-corporativo/files/content/file/2025/02/21/104/el-acuerdo-de-marco-avanzado-entre-la-ue-y-chile.pdf}{\textit{El Acuerdo de Marco Avanzado entre la UE y Chile}. CEOE, 2025}

\href{https://comercio.gob.es/PoliticaComercialUE/AcuerdosComerciales/Paginas/Chile.aspx}{\textit{Chile - Acuerdos comerciales de la UE}. Ministerio de Economía, Comercio y Empresa, 2026}

\href{https://www.subrei.gob.cl/sala-de-prensa/noticias/detalle-noticias/2024/11/13/congreso-aprueba-y-despacha-acuerdo-marco-avanzado-entre-chile---uni%C3%B3n-europea}{\textit{Congreso aprueba y despacha Acuerdo Marco Avanzado entre Chile y Unión Europea}. SUBREI, 2024}

\href{https://trade.ec.europa.eu/access-to-markets/es/content/acuerdo-comercial-interino-ue-chile}{\textit{Acuerdo comercial interino UE-Chile}. Access2Markets, Comisión Europea, 2025}

\href{https://obtienearchivo.bcn.cl/obtienearchivo?id=repositorio%2F10221%2F36448%2F1%2FBCN_Informe_AMA_Chile_UE.pdf}{\textit{Acuerdo Marco Avanzado Chile – Unión Europea}. Biblioteca del Congreso Nacional de Chile, 2025}

\href{https://isds-americalatina.org/chile/}{\textit{Chile}. ISDS América Latina, 2024}

\href{https://www.subrei.gob.cl/acuerdos-comerciales/modernizacion-ue-chile}{\textit{Acuerdo Marco Avanzado entre Chile y la Unión Europea}. SUBREI, s.f.}

\href{https://embamex.sre.gob.mx/belgica/index.php/es/avisos/670-agm}{\textit{Acuerdo Global Modernizado México-UE}. Embajada de México en Bélgica, 2026}

\href{https://trade.ec.europa.eu/access-to-markets/es/content/acuerdo-global-ue-mexico}{\textit{Acuerdo Global UE-México}. Access2Markets, Comisión Europea, 2026}

\href{https://imco.org.mx/el-acuerdo-global-modernizado/}{\textit{El Acuerdo Global Modernizado}. Instituto Mexicano para la Competitividad (IMCO), 2026}

\href{https://www.pan.senado.gob.mx/2026/07/celebra-marko-cortes-ratificacion-de-acuerdos-global-modernizado-y-comercial-interino-pero-senala-puntos-criticos-en-mexico-que-preocupan-a-la-ue/}{\textit{Celebra Marko Cortés ratificación de Acuerdos Global Modernizado y Comercial Interino}. Grupo Parlamentario del PAN, Senado de México, 2026}

\href{https://expansion.mx/economia/2026/07/08/parlamento-europeo-aprueba-acuerdo-modernizado-mexico}{\textit{Parlamento Europeo da luz verde al nuevo acuerdo comercial con México y lo aprueba con 479 votos}. Expansión, 2026}

\href{https://www.eeas.europa.eu/delegations/m%C3%A9xico/relaciones-ue-m%C3%A9xico-el-consejo-refrenda-los-acuerdos-para-impulsar-la-cooperaci%C3%B3n-y-el-comercio_es}{\textit{Relaciones UE-México: el Consejo refrenda los acuerdos para impulsar la cooperación y el comercio}. EEAS, 2026}

\href{https://www.gob.mx/sre/prensa/se-firma-el-acuerdo-global-modernizado-en-el-marco-de-la-viii-cumbre-mexico-union-europea}{\textit{Se firma el Acuerdo Global Modernizado en el marco de la VIII Cumbre México-Unión Europea}. Secretaría de Relaciones Exteriores de México, 2026}

\href{https://www.oxfamintermon.org/es/nota-de-prensa/60-organizaciones-exigen-suspension-acuerdo-asociacion-ue-israel}{\textit{Más de 60 organizaciones exigen la suspensión del Acuerdo de Asociación UE-Israel}. Oxfam Intermón, 2026}

\href{https://www.amnesty.org/es/latest/campaigns/2026/05/eu-israel-trade-agreement/}{\textit{Líneas rojas, no alfombras rojas: por qué Italia y Alemania deben apoyar la suspensión del Acuerdo de Asociación UE-Israel}. Amnistía Internacional, 2026}

\href{https://cde.ugr.es/index.php/union-europea/noticias-ue/2033-propuesta-de-suspension-de-concesiones-comerciales-con-israel-y-sanciones-a-los-ministros-extremistas-y-colonos-violentos}{\textit{Propuesta de suspensión de concesiones comerciales con Israel y sanciones a los ministros extremistas y colonos violentos}. Centro de Documentación Europea, Universidad de Granada, 2025}

\href{https://www.amnesty.org/es/latest/news/2026/04/eu-failure-to-suspend-eu-israel-association-agreement-shows-contempt-for-civilian-lives/}{\textit{UE: la decisión de no suspender el Acuerdo de Asociación de la Unión Europea e Israel denota desprecio por las vidas civiles}. Amnistía Internacional, 2026}

\href{https://fundacionalternativas.org/wp-content/uploads/2022/07/xmlimport-IbjEJY.pdf}{\textit{Los Acuerdos de Asociación Económica (EPAs) de la Unión Europea con África, el Caribe y el Pacífico}. Fundación Alternativas, s.f.}

\href{https://comercio.gob.es/PoliticaComercialUE/AcuerdosComerciales/Paginas/sudafrica.aspx}{\textit{Sudáfrica}. Ministerio de Economía, Comercio y Empresa, 2026}

\href{https://eur-lex.europa.eu/ES/legal-content/summary/trade-development-and-cooperation-agreement-tdca-with-south-africa.html?fromSummary=07}{\textit{Acuerdo de Comercio, Desarrollo y Cooperación (ACDC) con Sudáfrica}. EUR-Lex, 2018}

\href{https://ipsnoticias.net/2011/08/africa-austral-ue-asociacion-economica-extenuada/}{\textit{África Austral-UE: Asociación económica extenuada}. IPS Agencia de Noticias, 2011}

\href{https://www.bilaterals.org/?comercio-africa-union-aduanera-en&lang=en}{\textit{Comercio-África: Unión aduanera en riesgo ante los EPA}. bilaterals.org, s.f.}

\href{https://www.tralac.org/resources/our-resources/12127-sadc-eu-economic-partnership-agreement-documents-and-resources.html}{\textit{SADC-EU Economic Partnership Agreement Documents and Resources}. tralac Trade Law Centre, s.f.}

\href{https://trade.ec.europa.eu/access-to-markets/es/content/epa-sadc-southern-african-development-community}{\textit{AAE SADC - Comunidad para el Desarrollo del África Meridional}. Access2Markets, Comisión Europea, 2026}

\href{https://www.exteriores.gob.es/es/Comunicacion/NotasPrensa/Paginas/2023_NOTAS_P/20231115_NOTA042.aspx}{\textit{La Unión Europea y la Organización de Estados de África, el Caribe y el Pacífico firman el Acuerdo de Samoa}. Ministerio de Asuntos Exteriores de España, 2023}

\href{https://eur-lex.europa.eu/ES/legal-content/summary/partnership-agreement-between-the-eu-and-its-member-states-of-the-one-part-and-the-members-of-the-organisation-of-african-caribbean-and-pacific-states-of-the-other-part-the-samoa-agreement.html?fromSummary=11}{\textit{Acuerdo de asociación entre la Unión Europea y los miembros de la Organización de Estados de África, el Caribe y el Pacífico (Acuerdo de Samoa)}. EUR-Lex, 2023}

\href{https://www.consilium.europa.eu/es/press/press-releases/2023/11/15/samoa-agreement-eu-and-its-member-states-sign-new-partnership-agreement-with-the-members-of-the-organisation-of-the-african-caribbean-and-pacific-states/}{\textit{Acuerdo de Samoa: la UE y sus Estados miembros firman un nuevo acuerdo de asociación con los miembros de la OEACP}. Consejo de la Unión Europea, 2023}

\href{https://www.consilium.europa.eu/es/policies/cotonou-agreement/}{\textit{Acuerdo posterior a Cotonú}. Consejo de la Unión Europea, s.f.}

\href{https://ecpr.eu/Events/Event/PanelDetails/14584}{\textit{From Post-Cotonou to Samoa: The remaking of OACPS-EU relations}. European Consortium for Political Research (ECPR), 2024}

\href{https://es.wikipedia.org/wiki/Comunidad_Econ%C3%B3mica_de_Estados_de_%C3%81frica_Occidental}{\textit{Comunidad Económica de Estados de África Occidental}. Wikipedia, 2026}

\href{https://trade.ec.europa.eu/access-to-markets/es/content/epa-comunidad-de-africa-oriental}{\textit{EPA - Comunidad de África Oriental}. Access2Markets, Comisión Europea, s.f.}

\href{https://trade.ec.europa.eu/access-to-markets/es/content/epa-africa-occidental}{\textit{EPA - África Occidental}. Access2Markets, Comisión Europea, s.f.}

\href{https://origen-mercancias.es/es/acuerdos-preferenciales/aae-completos}{\textit{AAE Completos}. Origen de Mercancías, s.f.}

\href{https://trade.ec.europa.eu/access-to-markets/es/content/acuerdo-de-asociacion-economica-ue-cariforum}{\textit{Acuerdo de Asociación Económica UE-Cariforum}. Access2Markets, Comisión Europea, s.f.}

\href{https://www.southcentre.int/wp-content/uploads/2013/08/AN_EPA18_The-EU-CARIFORUM-EPA-on-Services-Investment-and-E-Commerce_ES.pdf}{\textit{Acuerdo de asociación económica CARIFORUM-CE: servicios, inversión y comercio electrónico}. South Centre, 2013}

\href{https://www.intracen.org/es/recursos/publicaciones/acuerdo-de-asociacion-economica-cariforum-ce-compromisos-asumidos-en-materia}{\textit{Acuerdo de asociación económica CARIFORUM-CE: compromisos asumidos en materia de servicios e inversiones}. Centro de Comercio Internacional (ITC), s.f.}

\href{https://www.wto.org/spanish/res_s/booksp_s/dispu_summary95_11_s.pdf}{\textit{Solución de diferencias en la OMC: Resúmenes de una página por caso, 1995-2011}. Organización Mundial del Comercio, 2012}

\href{https://www.fao.org/4/Y2732S/y2732s08.htm}{\textit{La situación de las preferencias comerciales en la OMC}. FAO, s.f.}

\href{https://www.wto.org/spanish/tratop_s/dispu_s/cases_s/ds27_s.htm}{\textit{Comunidades Europeas — Régimen para la importación, venta y distribución de bananos (DS27)}. Organización Mundial del Comercio, s.f.}

\href{https://www.ejil.org/article.php?article=535&issue=40}{\textit{The past, present and future ACP-EC trade regime and the WTO}. European Journal of International Law, s.f.}

\href{https://www.socialistsanddemocrats.eu/sites/default/files/2857_ES_EPA_Brochure_PSE_ES-web_6.pdf}{\textit{Acuerdos de Asociación Económica: el desarrollo como prioridad}. Grupo S\&D, Parlamento Europeo, s.f.}

\href{https://revistas.ucm.es/index.php/REDC/article/download/86897/4564456563263}{\textit{De Lomé a Cotonú}. Revista Electrónica de Estudios Internacionales, Universidad Complutense de Madrid, s.f.}

\href{https://www.asociacioneconomiacritica.org/wp-content/themes/vantage/JEC_old/jec13/Ponencias/economia%20mundial/LOS%20ACUERDOS%20DE%20ASOCIACION%20DE%20LA%20UNION.pdf}{\textit{Los Acuerdos de Asociación de la Unión Europea}. Asociación de Economía Crítica, 2013}

\href{https://eur-lex.europa.eu/legal-content/ES/TXT/HTML/?uri=CELEX%3A52021PC0579}{\textit{Propuesta de Reglamento por el que se aplica un sistema de preferencias arancelarias generalizadas}. Comisión Europea, EUR-Lex, 2021}

\href{https://www.boe.es/buscar/doc.php?id=DOUE-L-2023-81684}{\textit{Reglamento (UE) 2023/2663 del Parlamento Europeo y del Consejo, de 22 de noviembre de 2023}. BOE, 2023}

\href{https://www.iberley.es/legislacion/reglamento-delegado-ue-2025-1951-comision-29-septiembre-2025-modifican-anexos-ii-iv-vii-reglamento-ue-n-978-2012-parlamento-europeo-consejo-27285307}{\textit{Reglamento Delegado (UE) 2025/1951 de la Comisión, de 29 de septiembre de 2025}. Iberley, 2025}

\href{https://comercio.gob.es/ImportacionExportacion/Aranceles/Paginas/Sistema-de-preferencias-generalizado.aspx}{\textit{Sistema de preferencias generalizadas (SPG)}. Ministerio de Economía, Comercio y Empresa, 2026}

\href{https://www.taric.es/noticias/2023-12-21-novedades-spg/}{\textit{Novedades SPG}. Taric, 2023}

\href{https://trade.ec.europa.eu/access-to-markets/es/content/todo-menos-armas-eba}{\textit{Todo menos armas (EBA)}. Access2Markets, Comisión Europea, 2026}

\href{https://www.milkyfashions.com/bangladesh-ldc-graduation-2026/}{\textit{Bangladesh LDC Graduation 2026: What EU Buyers Must Know}. Milky Fashions, 2026}

\href{https://www.thedailystar.net/business/news/eus-gsp-lifeline-bangladesh-must-win-2029-4051411}{\textit{EU's GSP+: The lifeline Bangladesh must win before 2029}. The Daily Star, 2026}

\href{https://www.un.org/ldcportal/es/node/858}{\textit{Bangladesh graduation status}. Portal PMA, Naciones Unidas, 2026}

\href{https://www.thedailystar.net/business/economy/news/eu-gsp-rules-and-bangladeshs-rmg-4218431}{\textit{EU GSP rules and Bangladesh's RMG}. The Daily Star, 2026}

\href{https://cpd.org.bd/resources/2023/02/Graduating-from-the-LDC-Group-Challenges-Facing-Bangladesh.pdf}{\textit{Graduating from the LDC Group: Challenges Facing Bangladesh}. Centre for Policy Dialogue (CPD), Bangladesh, 2023}



\MyBibItem{DAWID y GATTI}{Chapter 2 - Agent-Based Macroeconomics}{2018}{https://doi.org/10.1016/bs.hescom.2018.02.006}


% ANEXOS
\clearpage
\anexo{\textbf{Preguntas Test}}
%\input{\TemaCode/questions.tex} 



\end{document}